\chapter{Primitive operators and geometric states}
\label{cyprim:chapter}

An automorphic product can prescribe the difference between even
and odd root multiplicities without prescribing either multiplicity.
The first interior root of the Igusa half-denominator is a useful
example. Its root space has dimensions $29\mid93$. An independently
defined elliptic coefficient on an Enriques surface has dimensions
$1\mid65$. Both have superdimension $-64$. A differential could remove
twenty-eight pairs, but it would also have to respect the real-root
operators. Their Serre relations detect an obstruction in height six.

The geometric operators carry further information. Incidence
correspondences on Hilbert schemes give creation and annihilation;
the surface diagonal gives joining and splitting. Elliptic
characteristic classes transport these operations to new lattices.
Their degeneration at $y=1$ has parameter torsion, higher products and
nontrivial derived fibres. Chiral de Rham sheaves provide states with
fields before an index is taken. Their normal-coordinate systems
make the remaining diagonal and Hall comparisons concrete.

We work over $\mathbb Q$ for ordinary cohomology and extend to
$\mathbb C$ for holomorphic fields. A complex is cohomologically
graded, with $C[r]^j=C^{j+r}$. For a complex of super vector spaces,
the sign in tensor interchange is the sum of cochain degree and
intrinsic parity. We specify a different grading whenever a
parameter is used to form a geometric cochain model.

\section{The isotropic face on an Enriques surface}
\label{cyprim:sec:face}

Let $T$ be a smooth projective complex Enriques surface. Its rational
cohomology is concentrated in degrees $0,2,4$, with dimensions
$1,10,1$. Choose an ample class $h$, put $d=\int_T h^2$, and normalize
$\omega\in H^4(T;\mathbb Q)$ by $\int_T\omega=1$. Then
\[
 H^*(T;\mathbb Q)=\mathbb Q1\oplus\mathbb Qh\oplus P\oplus\mathbb Q\omega,
 \qquad P=h^\perp\subset H^2(T;\mathbb Q),\qquad \dim P=9.
\]
The intersection pairing is nondegenerate and $d\ne0$.
These facts also follow from the canonical K3 cover
\cite[\S~3]{rewriteDolg2014}.

\begin{proposition}\label{cyprim:prop:face}
Make $P$ central and define brackets on the remaining three
directions by
\begin{equation}\label{cyprim:eq:face-bracket}
 [h,\omega]=-2\omega,\qquad [h,1]=2\cdot1,
 \qquad [\omega,1]=\frac2d h.
\end{equation}
This gives an even Lie algebra $\mathfrak l_T\cong\mathfrak{sl}_2\oplus P$
on $H^*(T;\mathbb Q)[2]$. With
\[
 E_+=\frac d2\omega,\qquad H_0=-h,\qquad F_-=-1,
\]
the adjoint action of $E_+$ is cup product by $h$.
The Lie subalgebra
\[
 \mathfrak b_T=\mathbb QF_-\oplus t\mathfrak l_T[t]
 \subset\mathfrak l_T[t]
\]
has root-space denominator
\begin{equation}\label{cyprim:eq:face-denominator}
 (1-y^{-1})\prod_{n\ge1}
       (1-p^ny)(1-p^n)^{10}(1-p^ny^{-1}),
\end{equation}
where $F_-$ has degree $(0,-1)$ and $t^nE_+,t^nH_0,t^nP,t^nF_-$
have degrees $(n,1),(n,0),(n,0),(n,-1)$, respectively.
\end{proposition}

\begin{proof}
The three displayed elements satisfy
\[
 [H_0,E_+]=2E_+,
 \qquad [H_0,F_-]=-2F_-,
 \qquad [E_+,F_-]=H_0.
\]
Together with central $P$, these relations prove Jacobi.
Moreover
\[
 [E_+,1]=h,\qquad [E_+,h]=d\omega,
 \qquad [E_+,\omega]=[E_+,P]=0,
\]
which is precisely $h\cup-$. Bracket adds the $t$-degree, so
$\mathfrak b_T$ is closed. The two elements $F_-,tE_+$ first
generate $tH_0,tF_-$ and then all three $\mathfrak{sl}_2$ directions
in every positive degree. Each space $t^nP$ is an additional
central summand. There are therefore ten directions at $(n,0)$
and one at each adjacent degree, giving
\eqref{cyprim:eq:face-denominator}.
\end{proof}

Both cubic relations
\[
 (\operatorname{ad}F_-)^3(tE_+)=0,
 \qquad (\operatorname{ad}(tE_+))^3F_-=0
\]
hold by the three-dimensional $\mathfrak{sl}_2$ calculation.
The nine central directions at each positive isotropic multiple
agree with the simple isotropic multiplicity in the
Gritsenko--Nikulin presentation \cite[Theorem~2.3]{GN97}.
This is an explicitly defined Lefschetz Lie structure. The
commutator of the cup product itself is zero, and the same
imaginary roots need not remain central after the third real
generator is included.

There are actual current fields for this algebra. Induce the
trivial module of $\mathfrak b_T\otimes\mathbb Q[z]$ to the loop
Lie algebra $\mathfrak b_T\otimes\mathbb Q[z,z^{-1}]$. Its vacuum
module has PBW space
\[
 U\bigl(\mathfrak b_T\otimes z^{-1}\mathbb Q[z^{-1}]\bigr),
\]
and the current of $a\in\mathfrak b_T$ is
$\sum_r(az^r)w^{-r-1}$. The loop bracket gives
$[a_\lambda b]=[a,b]$; translation sends $az^{-m}$ to
$m az^{-m-1}$. These are the universal level-zero current
fields of the face algebra. Their vacuum and their products
are defined algebraically; a compact Hall realization requires
an action comparison.

\section{Real Hall generators and primitive dimensions}
\label{cyprim:sec:real}

The spherical collection of Theorem~\ref{cycompact:thm:triple}
has Euler matrix
\begin{equation}\label{cyprim:eq:Cartan}
 G=\begin{pmatrix}2&-2&-2\\-2&2&-2\\-2&-2&2\end{pmatrix}
   =4I-2\mathbf1\mathbf1^{\mathsf t}.
\end{equation}
Let $Q$ be the quiver with two arrows $i\to j$ for each $i<j$.
Its doubled quiver has the preprojective algebra $\Pi_Q$,
obtained by imposing the moment relations
$\sum_{a:i\to j}a^*a-\sum_{a:j\to i}aa^*=0$ at each vertex,
with the rightmost map acting first. We use
finite-dimensional nilpotent $\Pi_Q$-modules: sufficiently
long paths act by zero.

For invariant constructible functions, define convolution by
Euler integration over submodules. The \emph{composition algebra}
$\mathcal C_Q$ is the subalgebra generated by the characteristic
functions $e_i$ of the vertex simples. Its unit is the
characteristic function of the zero module. Every iterated
product is the Euler integral over the corresponding module
flags, so forgetting intermediate subobjects proves associativity
in every finite arity.

Proposition~\ref{cycompact:prop:Serre} proves the cubic relations
directly from those flags. In dimension $(3,1)$ on two vertices,
the arrows give $X:\mathbb C^3\to\mathbb C^2$ and
$Y:\mathbb C^2\to\mathbb C^3$ with $YX=0$. A contributing
$b$-plane lies between $\operatorname{im}Y$ and $\ker X$.
If this inclusion is impossible the flag variety is empty;
otherwise its Euler characteristic is
$\binom{s-r}{b-r}$, with $r=\operatorname{rk}Y$ and
$s=3-\operatorname{rk}X$. The relation $YX=0$ implies
$\operatorname{rk}X+\operatorname{rk}Y\le2$, and hence
\[
 \sum_{b=0}^3(-1)^b\binom{s-r}{b-r}
       =(-1)^r(1-1)^{s-r}=0.
\]
After the divided-power factors are restored this is
$(\operatorname{ad}e_i)^3e_j=0$.

Lusztig's semicanonical realization theorem gives the stronger
identification
\begin{equation}\label{cyprim:eq:composition}
 U(\mathfrak n_G^+)\xrightarrow{\sim}\mathcal C_Q,
 \qquad e_i\longmapsto\mathbf1_{S_i},
\end{equation}
where $\mathfrak n_G^+$ is the positive Kac--Moody algebra
with matrix~\eqref{cyprim:eq:Cartan}. The precise nilpotent
constructible-function algebra and its flag basis are described
in \cite[\S\S~1.1,~2.1]{rewriteGLS2007}.
This theorem identifies the real composition algebra. It does
not identify all imaginary primitives of a cohomological Hall
algebra or transport their operations to compact K3 moduli.

The difference is visible in two exact quiver counts. Let
$A_{Q,\mathbf d}(v)$ count absolutely indecomposable
$Q$-representations of dimension $\mathbf d$ over a field with
$v$ elements.

\begin{proposition}\label{cyprim:prop:kac}
For the quiver $Q$,
\[
 A_{Q,(1,1,0)}(v)=v+1,
 \qquad
 A_{Q,(1,1,1)}(v)=(v+1)^2(v^2+2)
   =v^4+2v^3+3v^2+4v+2.
\]
\end{proposition}

\begin{proof}
At two vertices the two arrows form a nonzero vector in
$\mathbb F_v^2$, modulo scalar multiplication. This gives
$(v^2-1)/(v-1)=v+1$ classes. At three vertices each edge of
the triangle carries a two-vector. The representation is
indecomposable exactly when the nonzero edges connect the
vertices. Connectedness forces every endomorphism to be a
common scalar, over every field extension, so these objects
are absolutely indecomposable. Put $a=v^2-1$. There are
$3a^2+a^3$ connected assignments, and the effective vertex
torus acts freely with order $(v-1)^2$. Dividing gives the
asserted polynomial.
\end{proof}

For comparison with BPS spaces, retain their coefficient ring.
Let $\mathbb T$ be the arrow-scaling torus preserving the
canonical tripled-quiver potential, in the equivariant theory
of \cite[\S~1.4]{rewriteJN2026}.
For its full BPS space $\mathfrak g^{\mathbb T}_{\mathbf d}$,
the Kac--BPS identity has convention
\begin{equation}\label{cyprim:eq:BPS-dimension}
 \operatorname{grdim}\mathfrak g^{\mathbb T}_{\mathbf d}
 =(1-t)^{-\dim\mathbb T}A_{Q,\mathbf d}(t^{-1}),
 \qquad \operatorname{grdim}V=\sum_j\dim V^j\,t^{j/2}.
\end{equation}
The equivariant module is free. Reducing it along the
augmentation of $\mathbb Q[\operatorname{Lie}\mathbb T]$
therefore gives total dimensions $2$ and $12$ in the two
displayed charges \cite[\S~1.4, after~(19)]{rewriteJN2026}.
The current variable in Hall PBW is a separate factor.
Since $|\operatorname{sdim}V|\le\dim V$, these finite spaces
cannot be the Igusa root spaces of signed dimensions $10$
and $-64$. The conclusion concerns the finite parts in
\eqref{cyprim:eq:BPS-dimension}; restoring equivariant variables,
current modes or new geometric coefficients changes the comparison.

\section{The first interior root}
\label{cyprim:sec:root}

Let $\alpha_1,\alpha_2,\alpha_3$ have Gram matrix
\eqref{cyprim:eq:Cartan}. The eigenvalues are $4,4,-2$ and
the determinant is $-32$. Set
\[
 \delta_{ij}=\alpha_i+\alpha_j,
 \qquad \gamma=\alpha_1+\alpha_2+\alpha_3.
\]
The pair roots are isotropic and $(\gamma,\gamma)=-6$.
The Weyl vector $\rho=(\alpha_1+\alpha_2+\alpha_3)/2$
satisfies $(\rho,\alpha_i)=-1$ and $(\rho,\rho)=-3/2$.

The weak Jacobi form needed for these multiplicities is
\[
 \phi_{0,1}(q,y)
  =4\sum_{j=2}^4\frac{\vartheta_j(q,y)^2}{\vartheta_j(q,1)^2},
\]
where
\[
 \vartheta_2=\sum_{n\in\mathbb Z}q^{(n+1/2)^2/2}y^{n+1/2},
 \quad \vartheta_3=\sum_{n\in\mathbb Z}q^{n^2/2}y^n,
 \quad \vartheta_4=\sum_{n\in\mathbb Z}(-1)^nq^{n^2/2}y^n.
\]
The fractional powers cancel in the sum, giving
\begin{equation}\label{cyprim:eq:Jacobi-seed}
 \phi_{0,1}=y+10+y^{-1}
 +q(10y^2-64y+108-64y^{-1}+10y^{-2})+O(q^2).
\end{equation}
For example, the $q$ coefficient of the first ratio multiplied
by four is $2y^2-4+2y^{-2}$. The other two together contribute
$8((y+y^{-1})^2-8(y+y^{-1})+12)$.

Use the minimal Gritsenko--Nikulin presentation associated to
the Igusa half-denominator \cite[\S\S~2--3]{GN97}.
``Minimal'' here fixes the parity prescription: positive simple
supermultiplicities give even simple generators, negative ones
give odd simple generators, and no cancelling pairs are added.
There are three even real generators $e_i$ and nine even
imaginary simple generators at each $\delta_{ij}$. Write
$B_{ij}$ for the latter space. The full root exponent at
$\delta_{ij}$ is ten, including $[e_i,e_j]$; the exponent
at $\gamma$ is $-64$.

\begin{proposition}\label{cyprim:prop:root-parity}
The presentation has $93$ odd simple generators $c^u$ at
$\gamma$. Its root space at $\gamma$ has parity dimensions
\begin{equation}\label{cyprim:eq:root-parity}
 \dim\mathfrak g_\gamma=(29\mid93).
\end{equation}
\end{proposition}

\begin{proof}
Put $x_i=e^{-\alpha_i}$. Up to squarefree degree
$x_1x_2x_3$, the denominator after removing its Weyl monomial is
\[
 \prod_i(1-x_i)\prod_{i<j}(1-x_ix_j)^{10}
                    (1-x_1x_2x_3)^{-64}.
\]
Its coefficient of $x_1x_2x_3$ is $-1+30+64=93$.
No nonidentity term of the Weyl alternating sum contributes
there. A reduced word of length two already displaces the
Weyl vector by $3\alpha_i+\alpha_j$, of height four;
longer reduced words add positive roots. A simple reflection
has displacement $\alpha_i$. It fixes a null pair containing
$i$, so the resulting monomial repeats $\alpha_i$; it adds
$4\alpha_i$ to the opposite pair. Two imaginary corrections
have height at least four. The coefficient $93$ is therefore
the negative of the simple supermultiplicity at $\gamma$.
The minimal prescription makes these $93$ generators odd.

Real Serre relations start in height four. The multilinear
part of the free Lie algebra on $e_1,e_2,e_3$ has dimension two.
For each pair $\{i,j\}$, the nine brackets $[e_k,B_{ij}]$,
with $k$ the remaining index, give independent even vectors.
The relations $[e_i,B_{ij}]=[e_j,B_{ij}]=0$ occur in other
root degrees, and no further defining relation occurs here.
Thus the even dimension is $2+3\cdot9=29$. The odd part
consists of the $93$ simple generators.
\end{proof}

This calculation distinguishes simple multiplicities from
root multiplicities. It also distinguishes either parity
dimension from the signed exponent $29-93=-64$.
A finite Lie algebra retaining the calculation is obtained
from the generators $e_i,B_{ij},c^u$ by imposing the incident
commutators above and killing all root components of height
at least four. Call it $L_{\le3}$. Its height-one dimension
is $3$, its height-two dimension is $30$, and its height-three
dimensions are $35\mid93$: eight real Lie words and
twenty-seven mixed even brackets. Hence $\dim L_{\le3}=161$.
It is a Lie quotient, so Jacobi holds without a choice of
bracketing table.

Height truncation of a Lie algebra and height truncation of
its enveloping algebra are different operations. Even for
one primitive element $x$, the quotient $\mathbb C[x]/(x^{H+1})$
is not a Hopf quotient in characteristic zero: the coproduct
of $x^{H+1}$ has nonzero terms
$\binom{H+1}{j}x^j\otimes x^{H+1-j}$ with $0<j<H+1$.
By contrast, the enveloping algebra of a finite-height Lie
quotient is a Hopf algebra, generally infinite-dimensional.

\section{The canonical cover and an elliptic coefficient}
\label{cyprim:sec:elliptic-coefficient}

The canonical cover explains the half-index geometrically,
but its two characters must be retained when morphisms are
composed. Fix $K_T^{\otimes2}\cong\mathcal O_T$ and write
\[
 \pi:S=\operatorname{Spec}_T(\mathcal O_T\oplus K_T)\longrightarrow T.
\]
This is the \'etale K3 double cover. For perfect complexes
$F,G$ on $T$, put
\[
 \mathcal C^\epsilon(F,G)
 =R\operatorname{Hom}_T(F,G\otimes K_T^\epsilon),
 \qquad \epsilon\in\mathbb Z/2.
\]
Compose $f$ of character $\epsilon$ and $g$ of character
$\eta$ by $(g\otimes1_{K_T^\epsilon})f$, using the chosen
square trivialization. The character is independent of
cochain parity.

\begin{proposition}\label{cyprim:prop:cover}
There is a composition-compatible equivalence
\[
 R\operatorname{Hom}_S(\pi^*F,\pi^*G)
 \simeq\mathcal C^0(F,G)\oplus\mathcal C^1(F,G).
\]
The K3 Serre trace pairs opposite characters. The invariant
character alone is not a $2$-Calabi--Yau category.
\end{proposition}

\begin{proof}
Finite \'etale adjunction and the projection formula give
$\pi_*\pi^*G=G\otimes(\mathcal O_T\oplus K_T)$.
The algebra multiplication on this direct sum gives precisely
the stated composition. The deck involution acts by $-1$
on the K3 volume form, so the trace is nonzero on the
anti-invariant character; two paired morphisms must have
opposite characters. For $F=G=\mathcal O_T$, the invariant
endomorphisms have cohomology $(1,0,0)$, while the $K_T$
character has cohomology $(0,0,1)$. Removing the latter
removes the Serre partner of the identity.
\end{proof}

There is a related geometric caution. The free involution
$(\iota,-1)$ on $S\times E$ preserves its holomorphic
three-form, so the quotient has trivial canonical bundle.
Translation by $a\in E$ descends to this quotient precisely
when $2a=0$, since it must commute with $z\mapsto-z$.
This quotient does not retain the full translation action
used in the common-torsor construction of
Section~\ref{cycompact:sec:torsor}.

For a smooth projective variety $Y$ of even complex dimension
$d$, define the elliptic tensor class
\begin{equation}\label{cyprim:eq:elliptic-tensor}
 \mathcal E_{q,y}(Y)=y^{-d/2}\bigotimes_{r\ge1}
 \left(\Lambda_{-yq^{r-1}}T_Y^*
 \otimes\Lambda_{-y^{-1}q^r}T_Y
 \otimes\operatorname{Sym}_{q^r}T_Y^*
 \otimes\operatorname{Sym}_{q^r}T_Y\right).
\end{equation}
An exterior power of degree $a$ has intrinsic parity $a$;
its coefficient retains a positive monomial weight. Thus
the minus signs in~\eqref{cyprim:eq:elliptic-tensor} record
supercharacters rather than negative vector spaces.
Every $(q^my^\ell)$ coefficient is a finite superbundle.
Its derived global sections form a perfect supercomplex
$V_{m,\ell}(Y)$, with index
\[
 e_Y(m,\ell)=\sum_{j,\epsilon}(-1)^{j+\epsilon}
                       \dim H^j(V_{m,\ell}(Y))_\epsilon.
\]

\begin{proposition}\label{cyprim:prop:elliptic-coefficient}
For an Enriques surface, the $qy$ coefficient of
\eqref{cyprim:eq:elliptic-tensor} is internally even and equals
\begin{equation}\label{cyprim:eq:coefficient-bundle}
 B_{11}=\operatorname{Sym}^2\Omega_T^1\oplus K_T
                      \oplus\Omega_T^1\oplus T_T.
\end{equation}
Its cohomology dimensions in degrees $0,1,2$ are $(0,65,1)$,
so its total parity dimensions are $1\mid65$. The $qy^2$
coefficient is $\Pi T_T$ and has ten even cohomology classes.
On a K3 surface the $qy$ coefficient has cohomology
dimensions $(1,130,1)$.
\end{proposition}

\begin{proof}
Extracting energy one in the tensor product gives
\[
 \Omega_T^1\otimes\Omega_T^1
   \oplus K_T\otimes\Omega_T^1\oplus K_T\otimes T_T.
\]
Use $\Omega_T^1\cong T_T\otimes K_T$ and
$K_T^{\otimes2}\cong\mathcal O_T$ to obtain
\eqref{cyprim:eq:coefficient-bundle}. Every summand is even.
At Laurent weight two the only summand is the odd bundle
$K_T\otimes\Omega_T^1\cong T_T$.

A K3 surface has no positive symmetric holomorphic tangent
or cotangent tensors, by the Bochner and holonomy argument
\cite[Corollary~8]{rewriteKob1980}. Pullback to the canonical
cover gives the same vanishing for the corresponding Enriques
bundles and their canonical twists. Serre duality therefore
concentrates the cohomology of
$\operatorname{Sym}^2\Omega_T^1,\Omega_T^1,T_T$ in degree one.
Rationally $c_1(T_T)=0$, $\int_Tc_2=12$, and
\[
 \operatorname{ch}(\operatorname{Sym}^2\Omega_T^1)=3-4c_2,
 \qquad \operatorname{td}(T_T)=1+c_2/12.
\]
Riemann--Roch gives their degree-one dimensions $45,10,10$.
The canonical line has only $H^2(K_T)=\mathbb C$.
This proves $(0,65,1)$; the odd parity of $\Pi T_T$ changes
its ten degree-one classes to total even parity.
For K3 the three degree-one dimensions are $90,20,20$,
and $K_S=\mathcal O_S$ contributes in degrees zero and two.
\end{proof}

The coefficient bundle is the symbol of a field sheaf, whose
filtration need not split. Section~\ref{cyprim:sec:fields}
will recover the same dimensions for the actual chiral
coefficient by a field-theoretic projector. Neither calculation
identifies it with the root space~\eqref{cyprim:eq:root-parity}.

\subsection{Cyclotomic coefficient states}

The coefficient complexes themselves give another state
construction. On the energy-$r$ coefficient let $\mu_n$
act by $\zeta^r$, for a primitive $n$th root $\zeta$.
The averaging projector $n^{-1}\sum_{j=0}^{n-1}\zeta^{jr}$
retains exactly the energies divisible by $n$.
Regrade $r=nm$ to energy $m$ and assign length $n$.
This gives the complex
\[
 \mathcal P_Y=\bigoplus_{n>0,m\ge0,\ell\in\mathbb Z}
       p^nq^my^\ell V_{nm,\ell}(Y),
 \qquad
 \mathcal F_Y^{\mathrm{cyc}}
       =\widehat{\operatorname{Sym}}_{\mathrm{super}}(\mathcal P_Y).
\]
The differential extends as a derivation; the scalar at
length zero is the vacuum. Use total cochain and intrinsic
parity in the symmetry, and complete coefficientwise in
nonnegative length and energy. At fixed $p^Nq^M$, only
$n\le N,m\le M$ occur, and each tensor coefficient has
finitely many Laurent weights. Thus this completion is
defined before cancellation in the index.

For a homogeneous vector $v$ and a coefficientwise
continuous dual vector $\lambda$, creation by multiplication
and annihilation by the corresponding derivation satisfy
$[D_\lambda,M_v]_{\mathrm{super}}=\lambda(v)$.
The index is
\begin{equation}\label{cyprim:eq:cyclotomic-index}
 \operatorname{sch}\mathcal F_Y^{\mathrm{cyc}}
 =\prod_{n>0,m\ge0,\ell}
       (1-p^nq^my^\ell)^{-e_Y(nm,\ell)}.
\end{equation}
Indeed each finite complex splits over the ground field
into cohomology and contractible pairs. A contractible
even--odd pair contributes one to the symmetric-power
index; an even class of weight $x$ contributes
$(1-x)^{-1}$ and an odd class contributes $1-x$.
The finite coefficientwise product proves the formula.
This constructs complexes and state operators from the
elliptic tensor coefficients. A geometric twisted-sector
sheaf is defined separately below. Equal indices do not
identify these complexes or their geometric operations.

\section{A differential tested by the Serre relations}
\label{cyprim:sec:Serre}

Let $W_\gamma$ be the even part of the root space, and let
$R_\gamma$ be its subspace of Lie words in the real generators.
There is a direct sum decomposition
\[
 W_\gamma=R_\gamma\oplus
       [e_3,B_{12}]\oplus[e_2,B_{13}]\oplus[e_1,B_{23}],
 \qquad \dim R_\gamma=2.
\]
Choose independent $w_1,\ldots,w_{28}\in W_\gamma$.
Put the even generators in cochain degree zero and the
odd generators $c^1,\ldots,c^{93}$ in degree $-1$.
Their former super parity is now the parity of cochain
degree, with no additional intrinsic grading.
In $L_{\le3}$ the assignments
\begin{equation}\label{cyprim:eq:rank28}
 dc^u=\begin{cases}w_u,&u\le28,\\0,&u>28,\end{cases}
 \qquad de_i=dB_{ij}=0
\end{equation}
define a dg Lie algebra. Its degree-$\gamma$ complex has
cohomology $\mathbb C^{65}$ in degree $-1$ and $\mathbb C$
in degree zero. After choices, it is quasi-isomorphic to
$R\Gamma(T,B_{11})[2]$. This is a comparison of complexes
in one root degree.

The full positive presentation has more relations. Let
$e_i,f_i,h_i$ be the Chevalley triple for the $i$th real
root in the full root algebra, and write
$E_i=\operatorname{ad}e_i$, $F_i=\operatorname{ad}f_i$ and
$H_i=\operatorname{ad}h_i$. On $W_\gamma$,
\[
 [E_i,F_i]=H_i,\qquad H_iw=-2w.
\]
The real--imaginary Serre relation is $E_i^3c^u=0$, since
$1-(\alpha_i,\gamma)=3$.

\begin{lemma}\label{cyprim:lem:lowering}
The map $w\mapsto(F_1w,F_2w,F_3w)$ is injective on
$W_\gamma$, and $F_i^2W_\gamma=0$ for every $i$.
\end{lemma}

\begin{proof}
Use $u=[e_1,[e_2,e_3]]$ and $v=[e_2,[e_1,e_3]]$ as a
basis of $R_\gamma$. Jacobi and $[f_i,e_i]=-h_i$ give
\[
\begin{array}{c|cc}
 &u&v\\\hline
 F_1&4[e_2,e_3]&2[e_2,e_3]\\
 F_2&2[e_1,e_3]&4[e_1,e_3]\\
 F_3&-2[e_1,e_2]&2[e_1,e_2].
\end{array}
\]
For $b\in B_{ij}$ and the remaining index $k$,
\[
 F_k[e_k,b]=4b,\qquad F_i[e_k,b]=F_j[e_k,b]=0.
\]
The $B_{ij}$ directions are independent of the real pair
brackets. The three lowering maps first detect the twenty-seven
mixed coefficients; their first two rows on $R_\gamma$ have
determinant $12$. This gives an injective $30\times29$ matrix.
Its $i$th component lies in the pair space on the other two
indices, which $F_i$ kills.
\end{proof}

\begin{lemma}\label{cyprim:lem:sl2}
If $[E,F]=H$, $[H,E]=2E$, $[H,F]=-2F$,
$Hw=-2w$ and $F^2w=0$, then
\begin{equation}\label{cyprim:eq:144}
 F^4E^3w=144Fw.
\end{equation}
\end{lemma}

\begin{proof}
The commutation relations imply
\[
 FE^aF^bw=E^aF^{b+1}w
       -a(-2-2b+a-1)E^{a-1}F^bw.
\]
Four applications to $E^3w$ give successively
$E^3Fw$, $6E^2Fw$, $36EFw$ and $144Fw$.
\end{proof}

\begin{theorem}\label{cyprim:thm:Serre-obstruction}
Every root-preserving odd derivation of the unchanged
positive Borcherds presentation which satisfies $de_i=0$
and $dc^u\in W_\gamma$ has $dc^u=0$ for all $u$.
The rank-$28$ differential~\eqref{cyprim:eq:rank28} extends
to the positive Lie quotients of height at most five,
but it does not descend at height six.
\end{theorem}

\begin{proof}
Differentiate $E_i^3c^u=0$ to obtain $E_i^3dc^u=0$.
Equation~\eqref{cyprim:eq:144} gives $F_idc^u=0$ for all
$i$, and Lemma~\ref{cyprim:lem:lowering} gives $dc^u=0$.
The negative generators only detect a relation in the
positive algebra; the derivation need not act on them.
Square-zero was not used in this obstruction.

For the finite quotients, extend~\eqref{cyprim:eq:rank28}
by zero on every other simple generator. The first odd
simple roots have height three. A relation containing
$c^u$ is either a real Serre relation, first appearing at
height six, or a commutator for orthogonal simple roots.
But $(\gamma,\alpha_i)=-2$ and
$(\gamma,\delta_{ij})=-4$, so no such orthogonality
relation has height at most five. The remaining relations
are closed or discarded by height. The differential
preserves height and squares to zero on generators.
At height six the relations just tested prevent descent.
\end{proof}

There is already a loss of real-root information in degree
$\gamma$. Every $28$-dimensional subspace $W'\subset W_\gamma$
meets $R_\gamma$ in dimension at least
$28+2-29=1$. If $W'$ consists of boundaries and the real
generators are closed, a nonzero real Lie word becomes zero
in cohomology. The real composition algebra
\eqref{cyprim:eq:composition} cannot then act faithfully
through that cohomology with its original generators.

\subsection{A free attachment changes the presentation}

One can construct an all-height dg Lie algebra with the
desired boundaries by changing the relations. Let
$\mathfrak a$ be the ordinary Lie algebra generated by
the real and null simple generators with their prescribed
relations, in degree zero. Form the graded free product
\begin{equation}\label{cyprim:eq:attachment}
 L=\mathfrak a*\operatorname{Lie}(c^1,\ldots,c^{93}),
 \qquad |c^u|=-1,
\end{equation}
with root $\gamma$ for each $c^u$, and use
\eqref{cyprim:eq:rank28}. No real--imaginary Serre relation
is imposed on the new generators.

The free-product universal property extends the differential
uniquely. It squares to zero on generators and hence everywhere.
Its degree-zero cycles are $\mathfrak a$, and its degree-zero
boundaries form exactly the Lie ideal generated by the $w_u$.
Consequently
\[
 H^0(L)=\mathfrak a/\langle w_1,\ldots,w_{28}\rangle_{\mathrm{Lie}}.
\]
Positive-height completion preserves this construction:
every fixed root component is finite and the differential
preserves it. A homotopy-invariant derived quotient can
instead be formed from a semifree resolution of
$\mathfrak a$ and chosen closed lifts of the $w_u$.
Both constructions retain the changed quotient. Merely
adjoining a generator with $dz=E_i^3c^u$ would not repair
the old presentation, since $d^2z=E_i^3w_u$ can be nonzero.

Thus a geometric primitive comparison has two different
possibilities: retain the ordinary $29\mid93$ root space,
or construct a larger derived presentation together with
its geometric differential, brackets and syzygies. Equality
of the signed coefficient does not choose between them.

\section{Incidence and the surface diagonal}
\label{cyprim:sec:incidence}

Let $Y$ be a K3 or Enriques surface, put $X_n=Y^{[n]}$,
and write $A=H^*(Y;\mathbb Q)$. The vacuum space is
$H^*(X_0)=\mathbb Q$. Three coefficient theories occur
naturally on the same incidence geometry: constructible
functions, ordinary cohomology and coherent sheaves.
Their pushforwards and adjoints have different normalizations.

First consider zero-dimensional coherent sheaves on $Y$.
For constructible functions on their isomorphism classes,
the Euler Hall product is
\[
 (h_1*h_2)(Q)=\int_{U\subset Q}h_1(U)h_2(Q/U)\,d\chi.
\]
Lengths are fixed in each integral; constructibility is
understood on finite-type strata. A cyclic quotient is a
surjection $\mathcal O_Y\twoheadrightarrow Q$, and the
length-$n$ cyclic quotients form $X_n$.

\begin{proposition}\label{cyprim:prop:cyclic-action}
Constructible functions on the cyclic quotients carry the
unital Hall action
\[
 (h\cdot v)(\mathcal O_Y\twoheadrightarrow Q)
 =\int_{U\subset Q}h(U)
       v(\mathcal O_Y\twoheadrightarrow Q/U)\,d\chi.
\]
All finite iterates are compatible with forgetting stages
of the extension flag.
\end{proposition}

\begin{proof}
The quotient $Q/U$ inherits a surjective framing. Subobjects
of a fixed length are parameterized by a projective Quot
variety, so the integral is defined. Both orders of applying
$h_1,h_2$ integrate
$h_1(U)h_2(W/U)v(\mathcal O_Y\twoheadrightarrow Q/W)$ over
$0\subset U\subset W\subset Q$. Functoriality of Euler
pushforward proves the action identity and its all-flag
version. The length-zero subobject gives the unit.
\end{proof}

In geometric notation, $Z\subset Z'$ gives the sequence
\[
 0\longrightarrow I_Z/I_{Z'}\longrightarrow\mathcal O_{Z'}
   \longrightarrow\mathcal O_Z\longrightarrow0.
\]
The framing belongs to the middle object. This orientation
is necessary for the quotient to remain cyclic.

For $r>0$ let $\mathfrak Q_{n,r}$ be the incidence cycle
in $X_{n+r}\times Y\times X_n$ consisting of
$(Z',x,Z)$ with $Z\subset Z'$ and $I_Z/I_{Z'}$ of length
$r$ supported at $x$. Its fundamental cycle is the closure
of the generic punctual incidence locus. With projections
$\pi_+,s,\pi_-$, define
\[
 C_r(a)v=(\pi_+)_*
   \bigl([\mathfrak Q_{n,r}]\cup s^*a\cup\pi_-^*v\bigr).
\]
Here cap product and Poincar\'e duality are implicit; all
pushforwards are proper. Write $\dagger$ for the adjoint
under ordinary Poincar\'e integration and set
$A_r(a)=(-1)^{r-1}C_r(a)^\dagger$.
The geometric Heisenberg theorem gives
\begin{equation}\label{cyprim:eq:Heisenberg}
 [A_r(a),C_s(b)]=r\delta_{rs}\int_Yab,
 \qquad [C_r(a),C_s(b)]=[A_r(a),A_s(b)]=0.
\end{equation}
It identifies
\begin{equation}\label{cyprim:eq:geometric-Fock}
 \operatorname{Sym}\left(\bigoplus_{r\ge1}p^rA\right)
 \xrightarrow{\sim}\bigoplus_{n\ge0}H^*(X_n;\mathbb Q),
 \qquad p_r(a)\longmapsto C_r(a)|0\rangle,
\end{equation}
with products of generators interpreted as successive
creations \cite{Nakajima97,rewriteLS2003}.
The degree of $C_r(a)$ is $\deg(a)+2(r-1)$.
Over distinct points the iterated length-one correspondence
orders the points in $n!$ ways. Since both sides have degree
zero, this proves
\begin{equation}\label{cyprim:eq:unit-state}
 \frac1{n!}C_1(1)^n|0\rangle=1_{X_n}.
\end{equation}

There are also coherent kernels. On the nested scheme
$Z_{n,n+1}$ let $\mathcal L$ be its universal length-one
quotient line. For $F\in\operatorname{Perf}(Y)$ and
$k\in\mathbb Z$, the kernel
\[
 R(\pi_+,\pi_-)_*
   (\mathcal O_{Z_{n,n+1}}\otimes\mathcal L^k\otimes s^*F)
 \in\operatorname{Perf}(X_{n+1}\times X_n)
\]
is bounded coherent by properness and perfect by smoothness
of the product. Composition uses derived convolution. The
right adjoint kernel is its dual, tensored by the canonical
bundle of the source and shifted by its dimension.
For $n=k=0$ and $F=\mathcal O_Y$, creation is
$\mathbb C\mapsto\mathcal O_Y$ and its right adjoint is
$R\Gamma(Y,-)$. Their composite is $\mathbb C$ for Enriques
and $\mathbb C\oplus\mathbb C[-2]$ for K3. Thus the passage
to cohomological adjoints includes the Todd class; it cannot
be made by identifying the two transposes.

\subsection{Frobenius multiplication and graph genus}

Choose a basis $u_i$ of $H^2(Y)$, put
$Q_{ij}=\int_Yu_iu_j$, write $(Q^{ij})=Q^{-1}$,
and normalize $\int_Y\omega=1$.
Then $u_iu_j=Q_{ij}\omega$, and positive-degree classes
kill $\omega$. Write $m$ for cup multiplication. The diagonal
Gysin map, adjoint to the cup
product for $\int_Y$, is
\begin{equation}\label{cyprim:eq:Frobenius-diagonal}
\begin{split}
 \Delta1&=1\otimes\omega+\omega\otimes1
                  +\sum_{i,j}Q^{ij}u_i\otimes u_j,\\
 \Delta u_i&=u_i\otimes\omega+\omega\otimes u_i,
 \qquad \Delta\omega=\omega\otimes\omega.
\end{split}
\end{equation}
Testing against basis pairs proves the formulas. In
particular
\begin{equation}\label{cyprim:eq:Frobenius-Euler}
 e_A=m\Delta1=(b_2(Y)+2)\omega=c_2(T_Y),
 \qquad m\Delta(a)=e_Aa.
\end{equation}
A split followed by a join therefore inserts the Euler
class, with coefficient $12$ or $24$.

For $\sigma\in\mathfrak S_N$, label each cycle of
$\sigma$ by a factor of $A$ and form
$\mathscr A_N=\bigoplus_\sigma A^{\otimes c(\sigma)}\sigma$.
Here $c(\sigma)$ counts cycles. On a joint orbit $O$,
write $c_\sigma(O)$ for the number of cycles of its restriction.
To multiply a $\sigma$-labelled element by a
$\tau$-labelled element, consider each joint orbit $O$ of
$\langle\sigma,\tau\rangle$. Set
\[
 g_O=\frac{|O|+2-c_\sigma(O)-c_\tau(O)
                         -c_{\sigma\tau}(O)}2.
\]
The three permutations $\sigma,\tau,(\sigma\tau)^{-1}$
define a connected branched cover of the sphere on $O$.
Riemann--Hurwitz identifies $g_O$ as its genus.
Let $a_O,b_O$ be the products of the two sets of input
labels on $O$, and let $\Delta_k$ be the iterated diagonal,
with $\Delta_1=\operatorname{id}$. Multiply these labels, insert $e_A^{g_O}$,
and apply the iterated diagonal to obtain the labels on
the $c_{\sigma\tau}(O)$ output cycles. Explicitly, the
factor of the product on $O$ is
\begin{equation}\label{cyprim:eq:graph-product}
 \Delta_{c_{\sigma\tau}(O)}(a_Ob_Oe_A^{g_O}).
\end{equation}

The Frobenius identities make this multiplication associative:
cutting a connected covering surface into cylinders, pairs
of pants and handles gives multiplication, diagonal and
$m\Delta$, respectively. Associativity, coassociativity and
the Frobenius projection formula identify different cuts;
\eqref{cyprim:eq:Frobenius-Euler} accounts for every handle.
This is the Frobenius permutation construction and its
orbifold realization \cite[\S2]{rewriteLS2003}
\cite[\S~3]{rewriteFG2003}. For $s=(12)\in\mathfrak S_2$
and $r=(123)\in\mathfrak S_3$,
\[
 (as)(bs)=\Delta(ab)\,\mathrm{id},\qquad
 (ar)(br)=abe_A r^2,
 \qquad (ar)(br^{-1})=\Delta_3(ab)\,\mathrm{id}.
\]
The middle formula has genus one. Genus at least two gives
zero because $e_A^2=0$.

To obtain the Hilbert-scheme product, multiply each product
by
\begin{equation}\label{cyprim:eq:Hilbert-sign}
 \varepsilon(\sigma,\tau)
 =(-1)^{(|\sigma|+|\tau|-|\sigma\tau|)/2},
 \qquad |\sigma|=N-c(\sigma).
\end{equation}
The exponent is integral by permutation parity. The cocycle
identity follows by cancelling the middle lengths in a
triple product, so associativity is preserved.
The invariant algebra is $H^*(X_N;\mathbb Q)$ when the
canonical class is numerically trivial, including both
surfaces used here \cite[Theorem~3.2]{rewriteLS2003}.
The sum over conjugates, with repetitions, corresponds to
the product of incidence creators. Its identity sector is
therefore $N!$ times the unit, as in
\eqref{cyprim:eq:unit-state}. Over $\mathbb C$ the sector
rescaling by $i^{|\sigma|}$ removes the sign twist.

\subsection{Joining has order two}

In a homogeneous basis $a_\alpha$ of $A$, write
$a_\alpha a_\beta=m_{\alpha\beta}^{\gamma}a_\gamma$
and $\Delta a_\gamma=\Delta_\gamma^{\alpha\beta}
a_\alpha\otimes a_\beta$. Let $\partial_{i,\alpha}$
differentiate the Fock variable $p_i(a_\alpha)$. Define
\begin{equation}\label{cyprim:eq:join-split}
\begin{split}
 J&=\frac12\sum_{i,j\ge1}ij\,m_{\alpha\beta}^{\gamma}
       p_{i+j}(a_\gamma)\partial_{i,\alpha}\partial_{j,\beta},\\
 S&=\frac12\sum_{i,j\ge1}(i+j)\Delta_\gamma^{\alpha\beta}
       p_i(a_\alpha)p_j(a_\beta)\partial_{i+j,\gamma},
 \qquad D_H=J-S.
\end{split}
\end{equation}
Each fixed length sees finitely many terms. There are $ij$
choices of endpoints for joining cycles of lengths $i,j$.
A cycle of length $i+j$ has $i+j$ choices for a specified
split, with the ordered sum counting each split twice.
These counts give $J+S$ for the orbifold transposition
sum and $J-S$ for the Hilbert sign. Under the incidence
isomorphism, $D_H$ is multiplication by
$-c_1(\mathcal O_Y^{[N]})$; the canonical-class term in
Lehn's formula vanishes \cite[Theorem~3.10]{rewriteLehn1999}.

\begin{proposition}\label{cyprim:prop:joining}
If $M_i(a)$ is multiplication by $p_i(a)$, then
\[
 [[D_H,M_i(a)],M_j(b)]=ijM_{i+j}(ab),
 \qquad [[[D_H,M_i(a)],M_j(b)],M_k(c)]=0.
\]
The operator preserves length, raises cohomological degree
by two and normalized Hodge charge by one. It satisfies
\[
 D_H^2p_2(a)=-p_2(e_Aa),\qquad
 D_H^{2N+1}|_{H^*(X_N)}=0.
\]
\end{proposition}

\begin{proof}
Only the second-order term $J$ contributes to the double
commutator. Its two derivative slots give the displayed
product. A further commutator is zero. A join raises age
by one, contributing cohomological degree two; a split
adds diagonal degree four and lowers age by one.
Both have degree two and Hodge bidegree $(1,1)$.
On a cycle of length $i$, a class in $H^{r,s}(Y)$ has
normalized charge $r+i-1-i=r-1$, so the charge increases
by one. Splitting $p_2(a)$ and then joining yields
$-p_2(m\Delta a)$, proving the square. Finally the degrees
of $H^*(X_N)$ range from zero to $4N$.
\end{proof}

For $a=1$ the square is nonzero. In particular $D_H$
is even and second order; it is neither a
derivation of the Fock product nor a square-zero operator.

\section{Elliptic transport and its integral image}
\label{cyprim:sec:transport}

Put $R=\mathbb Q(y)[[q]]$ and
\[
 \epsilon_n(q,y)=\operatorname{td}(T_{X_n})
       \operatorname{ch}(\mathcal E_{q,y}(X_n)).
\]
Its integral is the elliptic genus. Its cohomological
degree-zero part $r_n$ has constant $q$ coefficient
$y^{-n}(1-y)^{2n}$ and is a unit of $R$.
The remaining part $\nu_n=\epsilon_n-r_n$ has positive
degree and $\nu_n^{2n+1}=0$. Hence
\begin{equation}\label{cyprim:eq:elliptic-inverse}
 \epsilon_n^{-1}=r_n^{-1}
       \sum_{j=0}^{2n}(-r_n^{-1}\nu_n)^j.
\end{equation}
Multiplication by $\epsilon_n$ is an invertible, generally
inhomogeneous map $U_n$ of ordinary cohomology over $R$.

A cohomology kernel $K$ from $X_m$ to $X_n$ acts by
pullback, cup product and integration along the source.
Write $p_n,p_m$ for the two projections of $X_n\times X_m$.
Define its transport by
\begin{equation}\label{cyprim:eq:transport}
 \widetilde K=p_n^*\epsilon_n\cup K\cup p_m^*\epsilon_m^{-1}.
\end{equation}

\begin{proposition}\label{cyprim:prop:transport}
Transport preserves identity kernels and all finite
convolutions. On operators it is
$T_{\widetilde K}=U_nT_KU_m^{-1}$.
\end{proposition}

\begin{proof}
In a convolution the two classes on every middle factor
multiply to $\epsilon_j^{-1}\epsilon_j=1$. The projection
formula moves the two outer classes through the remaining
pushforward. This proves the operator and composition
identities, in every finite arity. On the diagonal the two
classes cancel, giving the identity kernel.
\end{proof}

Thus the transported creators and annihilators satisfy
\eqref{cyprim:eq:Heisenberg}. Cup multiplication by the
tautological first Chern class commutes with $U_n$, so
$\widetilde D_H=D_H$. Let $\langle\mathcal I|$ integrate
each length component. Completing in length, the vacuum
normalization gives
\begin{equation}\label{cyprim:eq:elliptic-matrix}
 \frac1{n!}\widetilde C_1(1)^n|0\rangle=\epsilon_n,
 \qquad
 \left\langle\mathcal I\middle|
 e^{p\widetilde C_1(1)}\middle|0\right\rangle
   =\sum_{n\ge0}p^n\operatorname{Ell}(X_n;q,y).
\end{equation}
This constructs a specified matrix coefficient. The
parameter $q$ here belongs to the coefficient ring; the
inverse in~\eqref{cyprim:eq:elliptic-inverse} is an inverse
in cohomology. It is not a tensor inverse of the elliptic
coefficient complex.

There is considerable freedom in this construction.
Any even invertible classes $\alpha_n$ with $\alpha_0=1$
transport kernels in the same way. Taking
$\alpha_n=1+c_n[\mathrm{pt}]$ for $n>0$ gives any chosen
coefficients $c_n$ after integration. A scalar series
therefore does not characterize a critical Hall action
on an independently defined state object.

For integral behaviour near $y=1$, use $z=\log y$,
$\mathcal R=\mathbb Q[[q,z]]$, and
$H_n=H^*(X_n;\mathbb Q)\otimes\mathcal R$.
The class $\epsilon_n(q,e^z)$ is regular and becomes
invertible after inverting $z$. Its multiplication map
is injective because $H_n$ is free and the localized map
is invertible. Define the image lattice
\[
 \mathcal M_n=\epsilon_n H_n\subset H_n.
\]
Multiplication identifies $H_n$ with $\mathcal M_n$
as abstract free modules. A transported kernel preserves
these lattices exactly when the original kernel preserves
$H_m,H_n$, since
\begin{equation}\label{cyprim:eq:lattice-action}
 \widetilde K(\epsilon_m v)=\epsilon_n K(v).
\end{equation}
All rational geometric kernels above are therefore regular
on the image lattices. The hypothesis is necessary:
$z^{-1}$ times an identity kernel does not preserve any
nonzero free $\mathcal R$-lattice of this form.

\subsection{The Euler factor at every energy}

For a smooth projective $d$-fold $X$, with formal tangent
Chern roots $x_i$, put
\[
 e_z(TX)=\prod_{i=1}^d(x_i-z)
       =\sum_{k=0}^d(-z)^{d-k}c_k(TX),
 \qquad s(w)=\frac{\sinh(w/2)}{w/2}.
\]
Here $z$ is an even coefficient parameter, $c_0=1$ and
$s(0)=1$. Use the elliptic tensor formula with
$y^{-d/2}=e^{-dz/2}$ also when $d$ is odd, and put
$\epsilon_X=\operatorname{td}(TX)\operatorname{ch}(\mathcal E_{q,e^z}(X))$.
It has the factorization
\begin{equation}\label{cyprim:eq:Euler-factor}
 \epsilon_X(q,e^z)=e_z(TX)U_X(q,z),
 \qquad U_X\in(H^*(X)\otimes\mathcal R)^\times,
 \qquad U_X(q,0)=1,
\end{equation}
where
\[
 U_X=\prod_i\left[\frac{s(x_i-z)}{s(x_i)}
 \prod_{r\ge1}
 \frac{(1-e^zq^re^{-x_i})(1-e^{-z}q^re^{x_i})}
      {(1-q^re^{-x_i})(1-q^re^{x_i})}\right].
\]
Indeed the one-root, energy-zero factor is
\[
 e^{-z/2}\frac{x(1-e^{z-x})}{1-e^{-x}}
 =(x-z)\frac{s(x-z)}{s(x)}.
\]
The positive-energy factors are the second product.
Every fixed $q$ coefficient uses finitely many factors;
positive cohomological degrees are nilpotent. All remaining
factors are units and become one at $z=0$, proving the
identity by the splitting principle.

\begin{theorem}\label{cyprim:thm:Euler-quotient}
Let $d>0$, let $\pi:\mathbb P(TX)\to X$ be the bundle
of tangent lines, and put $\lambda=c_1(\mathcal O(-1))$. Then
\begin{equation}\label{cyprim:eq:projective-quotient}
 \frac{H^*(X;\mathbb Q)\otimes\mathcal R}{(\epsilon_X)}
 \cong H^*(\mathbb P(TX);\mathbb Q)[[q]],
 \qquad z\longmapsto\lambda.
\end{equation}
The image ideals of $\epsilon_X$ and $e_z(TX)$ coincide.
Their two-term multiplication complexes in resolution degrees
$-1,0$ are isomorphic
by $U_X$ on the source and the identity on the target.
\end{theorem}

\begin{proof}
The unit factorization proves the ideal and complex
assertions. The tautological inclusion
$\mathcal O(-1)\subset\pi^*TX$ is a nowhere-zero section
of $\pi^*TX\otimes\mathcal O(1)$, so
$\sum_k(-\lambda)^{d-k}\pi^*c_k(TX)=0$.
The classes $1,\lambda,\ldots,\lambda^{d-1}$ restrict
to a basis on each projective fibre. Leray--Hirsch
identifies the quotient by this monic relation with the
projective-bundle cohomology. The class $\lambda$ is
nilpotent, so completion in $z$ leaves the quotient
unchanged; adjoining $q$ gives the result.
\end{proof}

Write $D=\dim H^*(X;\mathbb Q)$. Over the discrete
valuation ring $\mathbb Q[[z]]$, multiplication by
$e_z(TX)$ has determinant $(-z)^{dD}$: in a basis ordered
by degree, each positive Chern class is strictly triangular.
It is therefore injective. Its positive Smith exponents
are the Jordan lengths of $\lambda\cup-$ on
$H^*(\mathbb P(TX))$, and their sum is $dD$.
The assertion extends to $\mathcal R$ by the unit
factorization, without invoking Smith normal form over a
two-variable ring. If $X$ is connected and $\chi(X)\ne0$,
multiplication at $z=0$ is by $c_d(TX)=\chi(X)[\mathrm{pt}]$
and has rank one. There is one unit Smith factor and
$D-1$ positive ones. The derived fibre of the cokernel
has rank $D-1$ in each resolution degree $-1,0$.
At length zero, $\epsilon_{X_0}=1$ and this cokernel is zero.

For a K3 or Enriques surface, $e_z(TY)=z^2+\chi(Y)\omega$.
On $\langle1,\omega\rangle$ its matrix is
$\left(\begin{smallmatrix}z^2&0\\\chi(Y)&z^2\end{smallmatrix}\right)$,
whose elementary divisors are $1,z^4$. Each of the
$b_2(Y)$ middle directions contributes $z^2$. Hence
\begin{equation}\label{cyprim:eq:surface-Smith}
 \operatorname{coker}\epsilon_Y
 \cong\mathcal R/(z^4)
       \oplus\bigl(\mathcal R/(z^2)\bigr)^{b_2(Y)}.
\end{equation}
The derived fibre has dimensions $11,11$ for Enriques
and $23,23$ for K3, over $\mathbb Q[[q]]$.
The abstract image lattice has rank $b_2(Y)+2$, while
its image in ambient cohomology at $z=0$ has rank one.
The difference is accounted for by the Tor term in the
specialization of $0\to\mathcal M_1\to H_1\to
\operatorname{coker}\epsilon_Y\to0$.

There is an explicit pole before passing to the image
lattice. At $q=0$, set $t=y-1$,
$s=y+y^{-1}-2=t^2/(1+t)$ and $a=1+s/12$.
The surface class is
\[
 \epsilon_Y(0,y)=s+a\,c_2(TY),\qquad
 [q]\epsilon_Y=-2s^2+
              \left(\frac56s^2-2s\right)c_2(TY).
\]
These follow by expanding the opposite Chern roots to
degree four. On $\langle1,\omega\rangle$ the inverse
matrix has lower-left entry $-a\chi(Y)s^{-2}$.
Thus $\widetilde A_1(1)(1)=\int_Y\epsilon_Y^{-1}$
has a fourth-order pole in $t$, although the elliptic
genus $\int_Y\epsilon_Y=a\chi(Y)$ is regular.

\subsection{The Hilbert square}

For $X=S^{[2]}$ with $S$ K3, the Fock decomposition gives
even Betti numbers $(1,23,276,23,1)$ and no odd cohomology.
Its Chern numbers and Chern pairing, for $\alpha,\beta\in H^2(X)$, are
\begin{equation}\label{cyprim:eq:Hilbert-square-Chern}
 \int_Xc_4=324,\qquad
 \int_Xc_2\alpha\beta=30q_X(\alpha,\beta),\qquad
 \int_Xc_2^2=828,
\end{equation}
where $q_X$ has its standard integral Beauville--Bogomolov
normalization. The first two are the Chern--Fujiki
identities \cite[Theorem~4.2]{rewriteCOT2023}; the last
also follows from $\chi(\mathcal O_X)=3$ and
$720\chi(\mathcal O_X)=3\int c_2^2-\int c_4$.
In particular $c_2:H^2(X)\to H^6(X)$ is an isomorphism.

\begin{proposition}\label{cyprim:prop:Hilbert-square-Smith}
The Smith exponents of Euler multiplication on $H^*(X)[[z]]$
are
\[
 0^1,\qquad2^{23},\qquad4^{276},\qquad6^{23},\qquad8^1.
\]
The same elementary divisors describe $\epsilon_X$ at
every elliptic energy. Its cokernel has derived fibre of
rank $323$ in each resolution degree $-1,0$.
\end{proposition}

\begin{proof}
Symplecticity gives $e_z(TX)=z^4+z^2c_2+c_4$.
Normalize $\int_X\omega=1$ and choose a basis $a_i$ of $H^2(X)$.
On $\langle1,c_2,\omega\rangle$ its matrix is
\[
 \begin{pmatrix}z^4&0&0\\z^2&z^4&0\\324&828z^2&z^4\end{pmatrix}.
\]
There is a unit entry. The two-by-two minors generate
$(z^4)$, since all are divisible by $z^4$ and the lower
left minor is $504z^4$. The determinant is $z^{12}$,
giving exponents $0,4,8$.
Each of the $23$ pairs $\langle a_i,c_2a_i\rangle$ in
degrees two and six has matrix
$\left(\begin{smallmatrix}z^4&0\\z^2&z^4\end{smallmatrix}\right)$
and exponents $2,6$. The kernel of $\int c_2(-)$ in
$H^4$ has dimension $275$ and is complementary to
$\mathbb Qc_2$; every such direction contributes $4$.
There are $323$ positive exponents, summing to
$1296=4\cdot324$. The unit factorization and the
two-term free resolutions prove the remaining assertions.
\end{proof}

These lengths constrain operators on the degeneration.
For $L_0=\mathbb Q[[z]]^D$ and
$M=\operatorname{diag}(z^{\ell_i})L_0$ with $\ell_i\ge0$, an ambient
endomorphism preserves $M$ precisely when
\begin{equation}\label{cyprim:eq:simultaneous-order}
 T_{ij}\in z^{\max(0,\ell_i-\ell_j)}\mathbb Q[[z]].
\end{equation}
This follows by conjugating with the diagonal matrix.
The inequality
$\max(0,a-b)+\max(0,b-c)\ge\max(0,a-c)$ proves closure
under composition. For exponents $(0,m)$, $m>0$, the
order has basis $e_{11},e_{22},e_{12},z^me_{21}$.
Its abstract fibre is four-dimensional, with two
opposite arrows whose composites vanish. Its image
on $L_0/zL_0$ is three-dimensional because the lower
arrow acts by zero. Thus even the algebra of regular
operators has distinct abstract and ambient fibres.

There cannot be a regular Heisenberg action on the
unchanged ambient modules with the transported creators
and the same empty vacuum. Indeed
$\widetilde C_1(\omega)|0\rangle
=z^2u(q,z)\omega$ with $u$ a unit. Every regular
annihilator sends this vector into $z^2\mathcal R$,
whereas~\eqref{cyprim:eq:Heisenberg} requires its
commutator with $\widetilde A_1(1)$ on the vacuum to
be $1$. The reversed product is zero since there are
no negative lengths. Equation~\eqref{cyprim:eq:lattice-action}
nevertheless supplies the regular action on the image
lattices themselves.

\section{Multiplicative and parameter fibres}
\label{cyprim:sec:higher-fibres}

Two gradings describe different aspects of the Euler
degeneration. The Smith calculation placed $z$ in
resolution degree zero. A geometric cochain model
instead gives $z$ degree two. Its derived fibre is a
sphere bundle and can have nonzero Massey products.
The unit $U_X$ in~\eqref{cyprim:eq:Euler-factor} is
generally inhomogeneous in ordinary cohomology, so
the equality of completed ideals does not identify
these two graded constructions.

Let $X$ be a compact smooth manifold and let $E\to X$
be a complex bundle of rank $r\ge1$.
Use rational polynomial forms $A_X=A_{\mathrm{PL}}(X)$,
with chosen closed Chern representatives, and put
\begin{equation}\label{cyprim:eq:projective-model}
 \mathcal P_E=A_X[z]\otimes\Lambda(\eta),
 \quad |z|=2,\quad |\eta|=2r-1,
 \quad Dz=0,\quad D\eta=e_z(E).
\end{equation}
This is a commutative cochain model of $\mathbb P(E)$.
To construct the comparison, send $z$ to a representative
of $c_1(\mathcal O(-1))$ and $\eta$ to a primitive of
the top Chern relation. Its induced cohomology map is
an isomorphism: multiplication by the monic polynomial
$e_z(E)$ on $H^*(X)[z]$ is injective, and the projective
bundle theorem identifies the resulting quotient.
Over real de Rham forms, that primitive has the explicit
Chern--Simons expression
\[
 \Theta=r\int_0^1P_r(\nabla_1-\nabla_0,F_t,\ldots,F_t)\,dt.
\]
Here $P_r(F,\ldots,F)=c_r(\nabla)$, $\nabla_1$ is
the tensor connection on $\pi^*E\otimes L^*$, and
$\nabla_0$ is trivial on its tautological nowhere-zero
line; $L=\mathcal O(-1)$, $\nabla_t=(1-t)\nabla_0+t\nabla_1$
and $F_t$ is its curvature. Write $\lambda=c_1(\nabla L)$.
The Bianchi identity gives
$d\Theta=c_r(\nabla_1)-c_r(\nabla_0)=e_\lambda(E)$.
Changing a Chern representative by $D\sigma$ changes
the relation generator by $\eta\mapsto\eta+\sigma$.

\begin{proposition}\label{cyprim:prop:sphere-fibre}
The multiplicative derived fibre at $z=0$ has model
\[
 \mathcal S_E=A_X\otimes\Lambda(\eta'),
 \qquad |\eta'|=2r-1,\quad D\eta'=c_r(E).
\]
It is the rational cochain algebra of the unit sphere
bundle $S(E)$. Equivalently,
\[
 C^*(\mathbb P(E);\mathbb Q)
       \otimes^{\mathbb L}_{\mathbb Q[z]}\mathbb Q
 \simeq C^*(S(E);\mathbb Q),\qquad |z|=2.
\]
\end{proposition}

\begin{proof}
Resolve $\mathbb Q$ over $\mathbb Q[z]$ by a generator
$\nu$ of degree one with $D\nu=z$. For
$Q_E(z)=(e_z(E)-c_r(E))/z$, set
$\eta'=\eta-\nu Q_E(z)$. Then $D\eta'=c_r(E)$.
The remaining factor $\mathbb Q[z]\otimes\Lambda(\nu)$
contracts to its constants. Geometrically, the unit circle
bundle of the tautological line is $S(E)$, the homotopy
fibre of its classifying map to $BU(1)$. Its odd-sphere
relative model adjoins the oriented fibre generator
with differential $c_r(E)$, giving the asserted model.
\end{proof}

\subsection{A nonzero ternary multiplication}

Suppose $X$ is connected and smooth projective, of
complex dimension $d>0$, with no odd rational cohomology
and $\chi=\chi(X)\ne0$. Set $A=H^*(X;\mathbb Q)$,
$\int_X\omega=1$ and $\tau(a)=\int_Xa$, extended by
zero outside top degree. K\"ahler formality allows the
tangent-sphere model to be written
\[
 B=A\otimes\Lambda(\eta),
 \qquad |\eta|=2d-1,\quad d\eta=\chi\omega
\]
\cite[\S~6]{rewriteDGMS1975}. Let
$A'=\bigoplus_{0<j<2d}A^j$. Representatives for its
cohomology are
\[
 H=\mathbb Q1\oplus A'\oplus\eta A^{>0}.
\]
The projection $p$ kills $\eta,\omega$, inclusion $i$
is the displayed representative inclusion, and the
homotopy is $h(\omega)=\eta/\chi$, zero elsewhere.
Thus $dh+hd=1-ip$ and $h^2=hi=ph=0$.

\begin{proposition}\label{cyprim:prop:ternary-fibre}
There is a strictly unital minimal $A_\infty$ model
of $B$ on $H$. Its binary multiplication is $p(i(-)i(-))$,
and its only possible nonzero higher multiplication is
\begin{equation}\label{cyprim:eq:ternary-fibre}
 m_3(a,b,c)=\frac{\tau(bc)a-\tau(ab)c}{\chi}\,\eta,
 \qquad a,b,c\in A'.
\end{equation}
It vanishes on other homogeneous input patterns, and
$m_n=0$ for $n\ge4$.
\end{proposition}

\begin{proof}
Fix suspension $s$ of degree $-1$ and bar signs
$b_1(sa)=-s(da)$,
$b_2(sa,sb)=(-1)^{|a|}s(ab)$.
The binary product loses exactly $\tau(ab)\omega$.
Its primitive is $h(ab)=\tau(ab)\eta/\chi$.
The two placements of this primitive in arity three
give $p(a\,h(bc)-h(ab)c)$, which is the displayed formula.
The comparison to $B$ has first term $i$, suspended
second term $I_2(sa,sb)=s(\tau(ab)\eta/\chi)$ on $A'$
and all later terms zero. The binary bar equation is
$-s\,d(\tau(ab)\eta/\chi)+s(ab)=s(ab-\tau(ab)\omega)$.

Every nonzero internal $h$ edge in a transfer tree has
output proportional to $\eta$. A second $h$ above it
vanishes; two such outputs multiply to zero. Every
tree with at least four leaves has one of these forms.
The same argument kills higher comparison terms and
all terms with an $\eta$-marked input. Normalization
of the contraction gives strict units. These identities
prove the complete bar equation, with $I_1$ inducing
the identity on cohomology.
\end{proof}

If $\dim H^2(X)\ge2$, choose independent $a,c\in A^2$
and $b\in A^{2d-2}$ with $\tau(ab)=1$, $\tau(bc)=0$.
Then the triple Massey product is
\begin{equation}\label{cyprim:eq:Massey}
 \langle[a],[b],[c]\rangle=-[c\eta]/\chi\ne0.
\end{equation}
Its indeterminacy is zero because $H^{2d-1}(B)=0$.
This proves nonformality of the tangent-sphere cochains.
It applies to every positive-length $Y^{[n]}$: the
incidence decomposition gives no odd cohomology, retains
the surface degree-two classes, and gives positive
Euler characteristic from
$\prod_{m\ge1}(1-p^m)^{-\chi(Y)}$.
For a surface the cohomology degrees are $0,2,5,7$,
with dimensions $1,b_2,b_2,1$, and
\[
 m_3(u_i,u_j,u_k)
   =\frac{Q_{jk}u_i-Q_{ij}u_k}{\chi(Y)}\eta.
\]
For a hyperbolic pair $u^2=v^2=0$, $uv=\omega$,
the choice $a=u+v$, $c=u-v$ gives
$\langle a,a,c\rangle=-2c\eta/\chi(Y)$.

\subsection{Jordan lengths as higher module operations}

Return to resolution grading: let $R_0=\mathbb Q[[z]]$
with $z$ of degree zero, and let $C$ be a finite torsion
$R_0$-module. Write $N$ for multiplication by $z$.
For $E=\Lambda(\xi)=\operatorname{Ext}_{R_0}^*(\mathbb Q,\mathbb Q)$,
$|\xi|=1$, a dg model of
$R\operatorname{Hom}_{R_0}(\mathbb Q,C)$ is
\begin{equation}\label{cyprim:eq:parameter-module}
 M_C=C\otimes\Lambda(\xi),
 \qquad d(a)=\xi Na,\quad d(\xi a)=0.
\end{equation}
This uses the free resolution with differential $-z$.
An additional internal degree two for $z$ is compatible
with internal degree $-2$ for $\xi$.

For $C_m=R_0/(z^m)$, choose $e_0,\ldots,e_{m-1}$ with
$Ne_j=e_{j+1}$ and $Ne_{m-1}=0$.
Cohomology has $u$ in resolution degree zero and $v$
in degree one. Representatives and homotopy are
\[
 i(u)=e_{m-1},\quad i(v)=(-1)^{m-1}\xi e_0,
 \quad h(\xi e_j)=e_{j-1}\ (j>0),
\]
with $h=0$ elsewhere. The projection sends $e_{m-1}$ to
$u$ and $\xi e_0$ to $(-1)^{m-1}v$, and kills the other basis vectors.
The unique nonzero reduced minimal-module operation is
\begin{equation}\label{cyprim:eq:Jordan-operation}
 m_{m+1}(\xi,\ldots,\xi,u)=v.
\end{equation}
Indeed the term with $j$ algebra inputs is
$p\xi(-h\xi)^{j-1}i$. Projection first survives at
$j=m$; the chosen sign of $v$ makes its coefficient one.
Every operation on $v$ vanishes. The comparison components
on the reduced bar comodule are
\[
 I_j((s\xi)^{\otimes j},u)=(-1)^je_{m-1-j}\quad(0\le j<m),
 \qquad I_0(v)=(-1)^{m-1}\xi e_0.
\]
The equation $dI_j+\xi I_{j-1}
=\sum_{k=1}^j I_{j-k}m_{k+1}$ cancels for $j<m$,
equals $I_0(v)$ for $j=m$, and is zero above it.
The minimal coderivation squares to zero since its
output is $v$. This verifies every arity, not just
the first nonzero one.

Applied to $C=H^*(\mathbb P(TX))$ with $N=\lambda\cup-$,
the multiplicities of these operations are
\[
\begin{array}{c|rrrr}
 X&m_3&m_5&m_7&m_9\\\hline
 \text{Enriques surface}&10&1&0&0\\
 \text{K3 surface}&22&1&0&0\\
 \text{K3 Hilbert square}&23&276&23&1.
\end{array}
\]
These are higher module actions in resolution grading.
The ternary multiplication~\eqref{cyprim:eq:ternary-fibre}
belongs to the separate multiplicative sphere model.
Deleting the module action~\eqref{cyprim:eq:Jordan-operation}
loses the relation $z^m$ even though the two cohomology
dimensions remain one.

Maps commuting with $N$ induce strict maps of
\eqref{cyprim:eq:parameter-module}. After transfer,
write $F(T)=P'TI$ for bar projection, the strict map
and bar inclusion. If the bar homotopy $H$ satisfies
$IP-1=-dH-Hd$, then
\[
 F(T_2)F(T_1)-F(T_2T_1)=dK+Kd,
 \qquad K=-P_3T_2H_2T_1I_1.
\]
The perturbation sums are finite at fixed bar length.
For two blocks $T_m=R_0/(z^m)$ and $T_n=R_0/(z^n)$,
their free resolutions also give
\[
 \operatorname{Hom}_{R_0}(T_m,T_n)
   =z^{\max(n-m,0)}R_0/(z^n),
 \qquad
 \operatorname{Ext}^1_{R_0}(T_m,T_n)=R_0/(z^{\min(m,n)}),
\]
and no higher Ext. Both nonzero groups have dimension
$\min(m,n)$. A correspondence that does not commute
with $N$ does not define such a parameter-module map.

\subsection{Gysin maps and maps of cones}

Some geometric maps act on the Euler cochain models
without a transfer. Let $E,F$ be bundles of ranks
$r\ge1,s\ge0$ on the same base, with compatible split
Chern representatives. For
$j:\mathbb P(E)\hookrightarrow\mathbb P(E\oplus F)$,
write $e_E=e_z(E)$, $e_F=e_z(F)$ and use relation
generators $\eta_E,\eta_{EF}$. The normal bundle is
$L^*\otimes F$. Pullback and Gysin have representatives
\begin{equation}\label{cyprim:eq:split-Gysin}
\begin{aligned}
 j^*(a)&=a,&j^*(\eta_{EF}b)&=\eta_Ee_Fb,\\
 j_!(a)&=e_Fa,&j_!(\eta_Eb)&=\eta_{EF}b,
 \qquad a,b\in A_X[z].
\end{aligned}
\end{equation}
Pullback is a dg algebra map. Gysin has degree $2s$
and is a cochain map: on the odd generator both
differentials give $e_Ee_Fb-\eta_{EF}Db$.
Multiplication verifies
\[
 j_!(j^*v\cdot w)=v\cdot j_!(w),\qquad
 j^*j_!(w)=e_Fw.
\]
Successive split inclusions multiply the appropriate
summand by $e_Fe_G\cdots$, so every finite composition
is independent of regrouping.

These maps represent the geometric operations. The
normal Euler class gives $j_!(1)=e_F(\lambda)$, and
the projective bundle is generated over the base by
$\lambda$, so the projection formula determines the
cohomology map. At the cochain level, take the universal
pair over $BU(r)\times BU(s)$. Its rational Chern-class
model is even. A degree-$2s$ map out of the Euler
resolution must send $\eta_E$ to the odd degree
$2(r+s)-1$, which is zero in the target quotient.
Its value on $1$ therefore determines it, and lifting
the universal Gysin map gives~\eqref{cyprim:eq:split-Gysin}.
Derived base change along the bundle classifying map
gives the assertion on $X$. This argument concerns
split vector bundles. A compact extension deformation
complex also contains cross derived-Hom terms; its
comparison with these bundles and its orientation are
additional requirements.

For general degeneration maps, a compatible square
includes its homotopy. If $f:A\to B$ and $g:C\to D$
are cochain maps, set
\[
 \operatorname{Cone}(f)=B\oplus A[1],\qquad
 d_f=\begin{pmatrix}d_B&f\\0&-d_A\end{pmatrix}.
\]
A filtered degree-$r$ map of cones has form
$T=\left(\begin{smallmatrix}b&h\\0&(-1)^ra\end{smallmatrix}\right)$,
with $|a|=|b|=r$, $|h|=r-1$. The differential
$\partial T=d_gT-(-1)^rTd_f$ and matrix composition
make these maps a dg category. In degree zero,
closedness is precisely the cochain condition on
$a,b$ together with
\begin{equation}\label{cyprim:eq:cone-square}
 d_Dh+hd_A=bf-ga.
\end{equation}
The composite homotopy is $b_2h_1+h_2a_1$.
These assertions follow by multiplying the two
matrices; $d_f^2=d_g^2=0$ gives $\partial^2=0$
and the graded Leibniz identity.

At length zero take $f_0=1:R_0\to R_0$, and at
length one take Euler multiplication
$f_1=e_z(TY):H^*(Y)[[z]]\to H^*(Y)[[z]]$
in resolution grading. The ordinary creator
$b(1)=1_Y$ admits no closed square: it would require
$1_Y=e_z(TY)a(1)$, contrary to the nonzero unit in
$H^*(\mathbb P(TY))$. The transported creator
$b(1)=e_z(TY)$ admits $a(1)=1_Y,h=0$.
Its cone map is nevertheless zero in the derived
category because $\operatorname{Cone}(f_0)$ is
contractible. Regularity on an image lattice therefore
does not supply a nonzero vacuum action on its cone.

\section{Normal coordinates and derived morphisms}
\label{cyprim:sec:formal-fibre}

A compact sheaf on a formal fibre has more information
than a sheaf with zero normal operator. For a smooth
projective surface $Y$, coherent sheaves on
$Y\times\operatorname{Spf}\mathbb C[[t]]$ killed by
some $t^M$ are equivalent to pairs $(F,N_F)$ with
$N_F$ nilpotent. Restriction of scalars and the rule
$t\mapsto N_F$ are inverse. They preserve kernels,
cokernels, extensions and finite flags. On flat
families with a common nilpotence bound they commute
with base change.

Choose a dg enhancement $C(F,G)$ of the surface Hom
complexes with associative composition and strict
closed nilpotent representatives for $N_F$.
Put $\delta_{GF}(a)=N_Ga-aN_F$ and adjoin
$|\theta|=1$. A model of the formal-fibre Hom is
\begin{equation}\label{cyprim:eq:normal-Hom}
\begin{split}
 \mathcal K_N(F,G)&=C(F,G)\otimes\Lambda(\theta),\\
 D(a\theta^\epsilon)&=d_Ca\,\theta^\epsilon
       +(-1)^{|a|}\delta_{GF}(a)\theta^{\epsilon+1},\\
 (a\theta^\epsilon)(b\theta^\eta)
       &=(-1)^{\epsilon|b|}ab\theta^{\epsilon+\eta}.
\end{split}
\end{equation}
The resolution $[\mathcal O_Y[t]\otimes F
\xrightarrow{t-N_F}\mathcal O_Y[t]\otimes F]$
together with a surface resolution gives the two copies
of $C(F,G)$ and their commutator
map. The exterior diagonal gives composition.
Equivalently, this is the one-dimensional Lie algebra
cochain model for the action $\delta$. Since $d_C$
and $\delta$ commute and satisfy the Leibniz rule,
$D^2=0$ and the displayed composition is a dg product.
No formality of $C(F,G)$ is required.

For K3 with a chosen holomorphic two-form, write
$\tau_Y:C(F,F)\to\mathbb C[-2]$ for its Serre trace.
The threefold trace is its $\theta$ coefficient,
normalized by the normal one-form.
In particular the pairing of $a\theta^\epsilon$ and
$b\theta^\eta$ is $(-1)^{\epsilon|b|}\tau_Y(ab)$
if $\epsilon+\eta=1$, and is zero otherwise.
Cyclicity and the vanishing of a commutator trace
give compatibility with the differential.

Let $f(t)=\sum_{r\ge1}a_rt^r$ with $a_1\ne0$ be
a new formal coordinate. For composable morphisms
$b_i:F_{i-1}\to F_i$, define
\begin{equation}\label{cyprim:eq:divided-difference}
 D_n^f(b_n,\ldots,b_1)=\sum_{r\ge n}a_r
 \sum_{j_0+\cdots+j_n=r-n}
 N_{F_n}^{j_n}b_nN_{F_{n-1}}^{j_{n-1}}
       \cdots b_1N_{F_0}^{j_0}.
\end{equation}
At fixed arity, nilpotence makes this finite.
The first difference satisfies
$D_1^f(N_Gb-bN_F)=f(N_G)b-bf(N_F)$.
Thus the linear map is the identity on unmarked
morphisms and $D_1^f$ on their $\theta$ components.
Its product defect is generally nonzero.

\begin{theorem}\label{cyprim:thm:coordinate-functor}
Use suspension $s$ of degree $-1$ and
$b_1(sx)=-s(Dx)$, $b_2(sx,sy)=(-1)^{|x|}s(xy)$.
For $n\ge2$, set
\[
 (\mathcal F_f)_n(s(b_n\theta),\ldots,s(b_1\theta))
       =s(D_n^f(b_n,\ldots,b_1)\theta),
\]
and set this component to zero on every other input
pattern. Together with the linear map just described,
these give a strictly unital $A_\infty$ functor
$\mathcal K_N\to\mathcal K_{f(N)}$. For invertible
coordinates $f,g$, there is an exact bar identity
\[
 \mathcal F_g\mathcal F_f=\mathcal F_{g\circ f},
\]
where the second coordinate is evaluated at $f(N)$.
In particular $\mathcal F_{f^{-1}}$ is an inverse.
\end{theorem}

\begin{proof}
The $d_C$ terms cancel by the graded Leibniz rule,
because every $N_F$ is closed of degree zero.
Among the remaining bar terms, all-marked inputs
have zero adjacent product since $\theta^2=0$.
Two or more unmarked inputs kill every higher term,
leaving only the ordinary binary identity. With
one unmarked input, the commutator inserts
$N_{F_i}b_i-b_iN_{F_{i-1}}$ in a divided difference.
For a monomial $t^r$ the powers at successive positions
telescope; the endpoints are the adjacent products.
For example
\[
\begin{split}
 D_1^f(ab)-aD_1^f(b)&=D_2^f(N_Ha-aN_G,b),\\
 D_1^f(a)b-D_1^f(ab)&=D_2^f(a,N_Gb-bN_F).
\end{split}
\]
At an interior position the same cancellation has
common suspended sign
$(-1)^{1+\sum_{i=k}^n|b_i|}$ on its three terms.
This proves the graded functor equation in every
arity. Linearity and fixed-arity finiteness extend
it to formal series.

For composition, put the endomorphisms $N_{F_i}$ on
the diagonal of an upper-triangular block matrix
and the $b_i$ on its first superdiagonal. The outer
block of $f$ of this matrix is $D_n^f$. The paths
through $g(f(-))$ run over all ordered partitions
of the inputs, exactly as in bar composition.
Functional calculus gives the asserted identity.
Mixed-input higher terms vanish on both sides.
An unmarked unit kills every higher component,
proving strict unitality.
\end{proof}

For $f(t)=t+at^2$ the correction terminates at
$D_2^f(x,y)=axy$, with
$D_1^f(x)=x+a(N_Gx+xN_F)$. A cubic coordinate
has a third component. The inverse coordinate
can have arbitrarily many components, each finite
on a fixed string of nilpotent objects.

\subsection{The coordinate contribution to the determinant}

For a cyclic trace and a formal series $h$, cyclic
permutation of the $r$ terms of $D_1^{t^r}$ gives
\begin{equation}\label{cyprim:eq:coordinate-trace}
 \tau_Y\bigl(h(N)f'(N)^{-1}D_1^f(b)\bigr)
       =\tau_Y(h(N)b).
\end{equation}
The inverse exists since $f'(0)\ne0$ and $N$ is
nilpotent. This transports the normal volume form
$h(t)\,dt$ and its trace pairing on cohomology.
It does not choose negative-cyclic representatives.

On a finite perfect equivariant chart over a local
Artin $\mathbb C$-algebra $R$, put
\[
 \chi_Y(F,G)=\sum_j(-1)^j\operatorname{rk}_R C^j(F,G).
\]
The determinant contribution is
\begin{equation}\label{cyprim:eq:coordinate-det}
 d_f(F,G)=\prod_j
 \det(D_1^f|_{C^j(F,G)})^{(-1)^{j+1}}.
\end{equation}
The extra sign comes from the degree of $\theta$.
At a complex point $D_1^f=a_1\operatorname{id}+U$
with $U$ nilpotent, so $d_f(F,G)=a_1^{-\chi_Y(F,G)}$.
Over a nonreduced base the nilpotent determinant
need not vanish. For K3, take
$R=\mathbb C[\epsilon]/(\epsilon^2)$,
$F=\mathcal O_Y\otimes R$, $N=\epsilon\operatorname{id}$,
and $f(t)=t+at^2$. Then
\[
 D_1^f=1+2a\epsilon,\qquad
 d_f(F,F)=1-4a\epsilon,\qquad d_f(F,F)^{1/2}=1-2a\epsilon.
\]

On K3 self-Hom, Serre duality makes $\chi_Y(F,F)$
even: degrees $j,2-j$ pair, and the middle pairing
is skew and nondegenerate. Put
$\rho_f(F)=a_1^{\chi_Y(F,F)}d_f(F,F)$.
Its reduction is one, and the square root with
specified reduction is uniquely
\begin{equation}\label{cyprim:eq:coordinate-half}
 j_f(F)=a_1^{-\chi_Y(F,F)/2}
        \exp\bigl(\tfrac12\log\rho_f(F)\bigr).
\end{equation}
Logarithm and exponential are finite on the Artin
maximal ideal. The linear part of the coordinate
functor proves the determinant cocycle; uniqueness
proves the square-root cocycle.
On a strict upper-triangular extension chart
$0\to F\to E\to G\to0$ whose Hom filtration is
preserved by the endomorphisms, the two cross
determinants agree by Serre adjunction. Consequently
\[
 d_f(E,E)=d_f(F,F)d_f(G,G)d_f(F,G)^2,\qquad
 j_f(E)=j_f(F)j_f(G)d_f(F,G).
\]
The second identity follows from the first and the
specified reductions. All finite flag refinements
use the same diagonal and unordered cross blocks;
compatible inverse limits of Artin charts preserve
the identities.

This supplies the normal-coordinate contribution
to square-root sewing on these charts. A global
orientation and stabilizer characters are separate
data. On an elliptic factor, a fixed invariant
one-form permits coordinates with translation
transitions. After recentering at a fibre the
transition is the identity, so $D_1=1$, $D_{n>1}=0$
and all these determinant factors are one. Descent
of critical Hall coefficient functors still needs
its supported operation maps.

\section{States with chiral fields}
\label{cyprim:sec:fields}

The elliptic tensor class has a field-bearing realization.
On a coordinate chart of a smooth projective complex $d$-fold $Y$,
the chiral de Rham sheaf $\Omega_Y^{\mathrm{ch}}$ has
even fields $\gamma^i,\beta_i$ and odd fields $c^i,b_i$,
with singular products
\begin{equation}\label{cyprim:eq:CDR-OPE}
 \beta_i(z)\gamma^j(w)\sim\frac{\delta_i^j}{z-w},
 \qquad b_i(z)c^j(w)\sim\frac{\delta_i^j}{z-w}.
\end{equation}
Their weights are $0,1,0,1$; the charges of $c,b$
are $1,-1$ and those of $\gamma,\beta$ are zero.
Functions of $\gamma$, translation and normal products
generate the local vertex superalgebra.
Write $\partial$ for vertex translation and $a_{(n)}b$
for its products; for $n\ge0$ these are the coefficients
of $(z-w)^{-n-1}$ in the singular OPE.
Let $L_0,J_0$ denote conformal weight and charge.
For $\widetilde\gamma=f(\gamma)$ with inverse $g$,
put $f_i^a=\partial_i f^a$ and $g_a^i=\partial_a g^i$.
The transformations are
\begin{equation}\label{cyprim:eq:CDR-coordinate}
\begin{split}
 \widetilde c^a&=f_i^a(\gamma)c^i,
 \qquad\widetilde b_a=:g_a^i(f(\gamma))b_i:,\\
 \widetilde\beta_a
   &=:\beta_i g_a^i(f(\gamma)):
     +:\partial_k(g_a^i\circ f)(\gamma)c^k b_i:.
\end{split}
\end{equation}
The order of the normal product in the last line
is part of the formula. The coordinate theorem
proves preservation of OPEs and composition on overlaps
\cite[\S~3]{rewriteMSV1999}.
For $\widetilde\gamma=\gamma^2$ on a punctured
one-dimensional chart, the last expression is
$:\beta(2\gamma)^{-1}:
-(2\gamma^2)^{-1}:cb:$; discarding the fermionic
term would change the coordinate map.

Choose a finite affine cover with affine intersections,
or a Stein cover acyclic for each coherent graded
piece. Let $V^\bullet$ be its cosimplicial section
object. The Thom--Whitney totalization
\[
 \mathbb V_Y=\operatorname{TW}(V^\bullet)
\]
uses compatible polynomial simplex forms tensored
with local sections. A form-field tensor has total
parity $\deg\omega+p(v)$, and its products are
\begin{equation}\label{cyprim:eq:TW-products}
 (\omega\otimes v)_{(n)}(\eta\otimes w)
 =(-1)^{p(v)\deg\eta}\omega\eta\otimes v_{(n)}w.
\end{equation}
The simplex differential is a derivation for these
signs. Translation acts on fields and the compatible
constant unit is the vacuum. Thus $\mathbb V_Y$
is a dg vertex superalgebra. Integration of simplex
forms compares it with the \v Cech complex; the
polynomial Poincar\'e contraction proves that this
comparison is a quasi-isomorphism. At fixed weight,
the coordinate construction has a finite filtration
by coherent tensor bundles. Acyclicity on the
intersections and induction on that filtration show
that its underlying complex computes
$R\Gamma(Y,\Omega_Y^{\mathrm{ch}})$ weight by weight.
Only the sheaf-resolution differential is included.

The associated-graded supercharacter is
\eqref{cyprim:eq:elliptic-tensor}. The filtration
need not split as a field sheaf. Nevertheless its
index is preserved, giving
\[
 \sum_i(-1)^i\operatorname{str}
 \bigl(q^{L_0}y^{J_0-d/2};H^i(Y,\Omega_Y^{\mathrm{ch}})\bigr)
 =\operatorname{Ell}(Y;q,y).
\]

\subsection{A projector onto the canonical descendant}

On the Enriques surface $T$, choose local square-root
volume forms for a fixed trivialization of $K_T^2$,
and volume coordinates for them. Their transition
determinants are constant $\pm1$. The transformation
of $J=\sum_i:c^ib_i:$ adds the vertex derivative of
the logarithm of that determinant, hence adds zero.
Thus $J$ is global. Its zero mode is the charge.
Set
$\mathscr V_{1,2}=(\Omega_T^{\mathrm{ch}})_{L_0=1,J_0=2}$;
this is the normalized $qy$ coefficient.
Weight zero and charge two are $K_T$.

\begin{proposition}\label{cyprim:prop:charge-projector}
As sheaves of complex vector spaces,
\[
 \mathscr V_{1,2}=\partial K_T\oplus\ker J_{(1)},
 \qquad P=\tfrac12\partial J_{(1)},\qquad P^2=P.
\]
The field coefficient has cohomology dimensions
$(0,65,1)$ in degrees zero, one and two.
The kernel has only degree-one cohomology, of
dimension $65$. For K3 the corresponding field
coefficient has dimensions $(1,130,1)$.
\end{proposition}

\begin{proof}
The translation relation is
$[J_{(1)},\partial]=J_{(0)}$. For $a\in K_T$,
weight gives $J_{(1)}a=0$ and charge gives
$J_{(0)}a=2a$. Hence $J_{(1)}\partial a=2a$,
proving the split injection and projector. These
maps need not be $\mathcal O_T$-linear.

The finite symbol filtration has factors
$\operatorname{Sym}^2\Omega_T^1,K_T,\Omega_T^1,T_T$.
Proposition~\ref{cyprim:prop:elliptic-coefficient}
therefore bounds the field cohomology by
$h^0=0$, $h^1\le65$, $h^2\le1$, and fixes its
Euler characteristic at $-64$. The split canonical
summand preserves $H^2(K_T)=\mathbb C$, so
$h^2=1$ and $h^1=65$. Removing that summand gives
the kernel assertion. On K3 the canonical summand
preserves one class in each of degrees zero and two;
the remaining symbol dimensions are $90,20,20$
in degree one, giving the stated result.
\end{proof}

This proves the actual total parity dimensions
$1\mid65$, independently of a splitting of the
symbol. It supplies no cancellation differential
on the $29\mid93$ Borcherds root space.

There is a natural additional differential, but it
has a different effect. Locally put
$Q=\sum_i:\beta_i c^i:$ and
$G=\sum_i:b_i\partial\gamma^i:$.
Their global operators satisfy
\begin{equation}\label{cyprim:eq:CDR-BRST}
 [Q_{(0)},G_{(1)}]_+=L_0.
\end{equation}
Coordinate changes of $Q$ differ by total vertex
derivatives, so its zero mode is global
\cite[\S~4]{rewriteMSV1999}.
Adding $Q_{(0)}$ to the resolution differential
contracts every positive-weight subcomplex, with
homotopy $G_{(1)}/m$ at weight $m>0$. The total
cochain signs make this homotopy anticommute with
the resolution differential. Moreover $Q_{(0)}$
raises charge by one. The state complex used here
therefore retains derived sections without this
extra differential.

\subsection{Permutation sectors and their operations}

The twisted chiral de Rham construction gives a
sheaf $\Omega_{Y^N}^{\mathrm{ch},g}$ on the fixed
locus of each $g\in\mathfrak S_N$, with a compatible
centralizer action \cite{rewriteFS2007}. Write $C(g)$
for the centralizer. The age is the sum of the fractions
$\theta\in(0,1)$ for its normal eigenvalues $e^{2\pi i\theta}$,
counted with multiplicity. Define
\begin{equation}\label{cyprim:eq:twisted-states}
 \mathcal H_N(Y)=\bigoplus_{[g]\subset\mathfrak S_N}
 R\Gamma((Y^N)^g,\Omega_{Y^N}^{\mathrm{ch},g})^{C(g)}.
\end{equation}
Use only the sheaf differential, twisted conformal
weight and charge $J_0-Nd/2$. Oscillator parity is
independent of the age contribution to $J_0$.
For an $n$-cycle, let
\[
 \mathcal K_n(Y)=
 R\Gamma(Y,\Omega_{Y^n}^{\mathrm{ch},(1\cdots n)})^{\mu_n}.
\]
Its oscillator indices lie in $n^{-1}\mathbb Z$.
For a surface, age $n-1$ makes the normalized
ground Laurent degree $-1$, as for the seed.

\begin{proposition}\label{cyprim:prop:cycle-decomposition}
Coefficientwise in length and conformal weight,
there is a natural equivalence of supercomplexes
\begin{equation}\label{cyprim:eq:cycle-Fock}
 \prod_{N\ge0}p^N\mathcal H_N(Y)
 \simeq\widehat{\operatorname{Sym}}_{\mathrm{super}}
       \left(\bigoplus_{n\ge1}p^n\mathcal K_n(Y)\right).
\end{equation}
\end{proposition}

\begin{proof}
A permutation with $a_n$ cycles of length $n$ has
fixed locus $Y^{\sum a_n}$ and centralizer
$\prod_n(\mu_n^{a_n}\rtimes\mathfrak S_{a_n})$.
The independent coordinate sets give external
products of the single-cycle twisted sheaves.
Their finite coherent filtrations justify
coefficientwise K\"unneth. Finite-group invariants
are exact over $\mathbb C$. Taking cyclic invariants
first and symmetric invariants second gives the
super symmetric powers, with cochain and intrinsic
parities combined as in~\eqref{cyprim:eq:TW-products}.
There are finitely many partitions at each length,
which proves the completed statement.
\end{proof}

On the symmetric model, concatenation of cycles
and assignment of each cycle to one of two groups
give a commutative, cocommutative dg super Hopf
algebra. In coinvariant coordinates,
\[
 \Delta(v_1\cdots v_r)=\sum_{I\subset\{1,\ldots,r\}}
 (-1)^{\epsilon(I)}v_I\otimes v_{I^c}.
\]
Order is retained in each subword and $\epsilon(I)$
is the total Koszul sign. In invariant coordinates
the multiplication is the averaged shuffle, fixing
the factorials. The empty cycle word is the unit,
single cycles are primitive and their antipode is
negation. Splitting a word into three groups proves
coassociativity directly. A continuous homogeneous
linear functional $\lambda$ on the single-cycle
space gives a derivation with
$[D_\lambda,M_v]_{\mathrm{super}}=\lambda(v)$.
All operations are finite at fixed length.

There is also an action of singular field modes.
For $V_Y=H^0(Y,\Omega_Y^{\mathrm{ch}})$, its mode
Lie superalgebra is
\[
 \mathfrak m(V_Y)=
 \frac{V_Y\otimes\mathbb C[t,t^{-1}]}
 {\operatorname{im}(\partial\otimes1+1\otimes\partial_t)},
 \qquad
 [a[m],b[n]]=
 \sum_{j\ge0}\binom mj(a_{(j)}b)[m+n-j].
\]
Here $a[m]$ is the class of $a\otimes t^m$.
The invariant field $a^\Delta=\sum_{i=1}^Na^{(i)}$
acts on every twisted sector by integer modes.
Write $\rho(a[m])=(a^\Delta)_{(m)}$ for this action.
For $j\ge0$, different factors have regular
products, so
$(a^\Delta)_{(j)}b^\Delta=(a_{(j)}b)^\Delta$.
Translation is preserved, and the central vacuum
mode $1[-1]$ acts by $N$. Under~\eqref{cyprim:eq:cycle-Fock}
the action is the sum on the individual cycles;
therefore
\begin{equation}\label{cyprim:eq:mode-creator}
 [\rho(A),M_v]_{\mathrm{super}}=M_{\rho(A)v}.
\end{equation}
The diagonal field assignment is not a unital
vertex map: its vacuum is multiplied by $N$,
and regular products have cross terms.

If $\operatorname{Ell}(Y;q,y)=\sum c(m,\ell)q^my^\ell$,
the twisted trace theorem and the symmetric-product
elliptic genus formula give
\begin{equation}\label{cyprim:eq:DMVV}
 \sum_{N\ge0}p^N\chi_s\mathcal H_N(Y)
 =\prod_{n>0,m\ge0,\ell}
       (1-p^nq^my^\ell)^{-c(nm,\ell)}
\end{equation}
\cite[Theorem~6.1]{rewriteFS2007}
\cite[Corollary~6.7]{rewriteBL2005}.
Here $\chi_s$ includes cochain and oscillator signs.
The same scalar series is the Hilbert-scheme elliptic
genus. Its equality does not construct an equivalence
of chiral coefficient complexes. The cycle product
has zero primitive commutator, while the seed mode
action preserves length; neither has yet been
identified with the compact extension action.

\section{An unshifted comparison detected by a field mode}
\label{cyprim:sec:McKay}

Let $S$ be a smooth projective K3 surface and set
$C_{H,N}=R\Gamma(S^{[N]},\Omega_{S^{[N]}}^{\mathrm{ch}})$
and $C_{O,N}=\mathcal H_N(S)$. Retain sheaf degree
$k$, weight $w$, charge $J$ and intrinsic parity
$\epsilon$. Total parity is $k+\epsilon$, and the
differential has degree $(1,0,0,0)$. A degree-zero
comparison preserving total parity, weight and
charge also preserves intrinsic parity at each
sheaf degree.

Every untwisted local monomial satisfies
$\epsilon\equiv J\pmod2$: each $c$ or $b$ contributes
one to parity and respectively $1$ or $-1$ to charge.
Coordinates, momenta, coefficient functions and
translation preserve this relation. The coordinate
maps do also. Thus the charge-even, intrinsically
odd part of $C_{H,N}$ is a zero sheaf complex, and
its cohomology is zero in every homogeneous model.

The field $G=\sum:b_i\partial\gamma^i:$ is global
on every smooth complex manifold. The antighosts
transform by the inverse transpose Jacobian and
coordinate derivatives by the Jacobian. All the
factors in these two transformations have regular
mutual OPEs, so the Jacobians cancel without an
extra contraction. The field is odd, of weight
two and charge $-1$.

In the transposition sector of $S^2$, let $\tau$
be the constant ground section. Its fixed locus
is $S$, and the natural centralizer action fixes
$\tau$. Its gradings are $(0,0,1,0)$: age contributes
one to charge but no oscillator parity. The global
invariant descendant
\begin{equation}\label{cyprim:eq:McKay-state}
 s=(G^{(1)}+G^{(2)})_{(0)}\tau
\end{equation}
has gradings $(0,1,0,1)$.

\begin{proposition}\label{cyprim:prop:McKay-detector}
The class of $s$ is nonzero in $H^0(C_{O,2})$.
Every degree-zero map $C_{H,2}\to C_{O,2}$ preserving
weight, charge and total parity has a nonzero
degree-zero comparison-cone class in this coefficient.
A reverse map has a nonzero cone class in degree
$-1$. There is no zigzag of quasi-isomorphisms
with these gradings.
\end{proposition}

\begin{proof}
Choose dual sum and difference coordinates locally.
Use oscillator expansions
$\gamma(z)=\sum_t\gamma_tz^{-t}$ and
$\beta(z)=\sum_r\beta_rz^{-r-1}$, with the same
exponents for $c$ and $b$, respectively.
The two normal coordinates have half-integral modes
with
$[\beta^-_{a,r},\gamma^{b,-}_t]
=\{b^-_{a,r},c^{b,-}_t\}
=\delta_a^b\delta_{r+t,0}$.
Positive normal modes annihilate $\tau$.
The normal part of the $G$ descendant is
\[
 s=\frac12\sum_{a=1}^2
       \gamma^{a,-}_{-1/2}b^-_{a,-1/2}\tau.
\]
The derivative supplies the factor $1/2$; the
integral tangential terms kill the constant ground
section. The actual mode operation
\begin{equation}\label{cyprim:eq:McKay-detector}
 2\beta^-_{1,1/2}c^{1,-}_{1/2}s=\tau
\end{equation}
proves nonvanishing. Globality was established
before the local coordinates were chosen. The
sheaf resolution has no negative degree, so a
nonzero global section gives a nonzero $H^0$ class.

For a forward map use
$\operatorname{Cone}(F)^k=C_{O,2}^k\oplus C_{H,2}^{k+1}$.
Its $(w,J,\epsilon)=(1,0,1)$ component is precisely
the orbifold component, since the Hilbert component
vanishes. It retains $[s]$. The reverse cone has
this component shifted by one, giving degree $-1$.
The same cohomology mismatch excludes a homogeneous
zigzag. Resolution changes, contractible additions
and coordinate changes preserve the nonzero class;
coefficientwise completion does not change its
fixed length and weight.
\end{proof}

The detector extends to every cycle length $N\ge2$.
Let $\tau_N$ be the constant ground section for
$(1\cdots N)$ and put $s_N=(\sum_iG^{(i)})_{(0)}\tau_N$.
In dual Fourier bases put $\theta_r=r/N$ for
$1\le r<N$. The coordinate and $c$ modes have
indices in $-\theta_r+\mathbb Z$, while their
momentum and antighost partners have indices in
$\theta_r+\mathbb Z$. The global state is
\[
 s_N=\sum_{a=1}^2\sum_{r=1}^{N-1}\theta_r
       \gamma^{a,r}_{-\theta_r}
       b_{a,r;-(1-\theta_r)}\tau_N,
\]
and each selected summand is detected by
\[
 \theta_r^{-1}\beta_{a,r;\theta_r}
            c^{a,r}_{1-\theta_r}s_N=\tau_N.
\]
The paired centralizer eigenvalues cancel. Age
$N-1$ gives gradings $(0,1,N-2,1)$ and normalized
charge $-2$. For every even $N$ the same untwisted
parity mismatch excludes the homogeneous comparison.
For odd $N$ this test gives no obstruction.

A different recipient can shift both gradings.
For a single $m$-cycle replace
\begin{equation}\label{cyprim:eq:age-shift}
 \mathcal K_m(S)\quad\text{by}\quad
       \Pi^{m-1}\mathcal K_m(S)[-(m-1)].
\end{equation}
This adds age to sheaf degree and intrinsic parity,
where $\Pi$ reverses intrinsic parity and the shifted
differential has sign $(-1)^{m-1}$. It preserves total
parity and the elliptic supertrace.
It removes the mismatch above: at length two $s$
moves to $(1,1,0,0)$. The shifted complex is not
degree-zero quasi-isomorphic to the original one,
whose ground class occurs in degree zero.
For a general permutation use age $N-c(g)$ in
each summand. Comparing this new recipient with
$C_{H,N}$ remains an open problem requiring field,
creation and joining maps with their supported
compatibilities. The grading change alone supplies
none of those maps.

\section{Resolving the target diagonal by normal fields}
\label{cyprim:sec:normal-diagonal}

Ordinary joining uses a diagonal Gysin map. A map
on fields has to retain the conjugate normal
variables as well. In one normal coordinate,
write the even fields as $v,\beta$ and the odd
fields as $c,b$, with
$\beta_{(0)}v=b_{(0)}c=1$.
A unital vertex homomorphism killing $v$ would
give $1=\phi(\beta)_{(0)}0=0$. Restriction of
functions therefore does not extend to such a
map on the full field system.

Instead define
\begin{equation}\label{cyprim:eq:normal-differential}
 \delta=(:vb:)_{(0)},
 \qquad \delta c=v,\quad\delta\beta=-b,\quad
 \delta v=\delta b=0.
\end{equation}
Use normal cochain degrees $|c|=-1$, $|b|=1$,
$|v|=|\beta|=0$. In this grading the normal
fermion parity is carried by cochain degree,
and is not counted a second time as intrinsic
parity. Tangential fermion parity remains
separate. The defining current has regular
self-OPE, so $\delta^2=0$. Its zero mode is a
derivation of every vertex product and commutes
with translation. The minus sign is forced by
$\delta(\beta_{(0)}c)=-b_{(0)}c+\beta_{(0)}v=0$.

For a formal chart, complete the coefficient
functions only in the undifferentiated even
normal coordinates at each fixed conformal
weight, and take the direct sum over weights.
Every state then has bounded conformal weight;
positive-weight oscillators remain polynomial.
For a pair of states the singular products
differentiate their coefficients only finitely
many times, so vertex products are continuous
on this carrier.

\begin{proposition}\label{cyprim:prop:normal-contraction}
On the polynomial carrier and on this formal
carrier, the inclusion of a tangential chiral
de Rham algebra into its tensor product with
finitely many normal quartets is a
quasi-isomorphism of dg vertex superalgebras.
\end{proposition}

\begin{proof}
In normalized derivative PBW variables,
\[
 \delta=\sum_j(v_j\partial_{c_j}-b_j\partial_{\beta_j}),
 \qquad
 h_0=\sum_j(c_j\partial_{v_j}-\beta_j\partial_{b_j}).
\]
Use left exterior derivatives. Their anticommutator
is the total normal-number operator $\mathcal N$,
since it is a derivation acting identically on
each generator. Sum over all normal directions.
On positive normal number put $h=h_0/\mathcal N$,
and put $h=0$ on the normal vacuum. If $i$ includes
the vacuum and $\pi$ extracts its coefficient,
\[
 \delta h+h\delta=1-i\pi.
\]
This is coefficientwise valid on formal functions;
differentiation lowers their adic order by at
most one. Tensor with the tangential factor.
Its inclusion preserves all products, translation
and vacuum, proving the asserted quasi-isomorphism.
\end{proof}

The projection $\pi$ is a complex map, not a
vertex map; $\beta_{(0)}v=1$ excludes the latter.
The homotopy is also only linear. Bounded weight
is essential: completing simultaneously to admit
$\sum_{N\ge0}:\beta^N:$ and $\sum_{N\ge0}:v^N:$
would admit complete contractions
$N!(z-w)^{-N}$ of unbounded pole order.
The differential preserves weight and lowers
the original fermionic charge by one. It is
therefore not a differential on each Laurent
coefficient separately. Unlike the extra de
Rham differential~\eqref{cyprim:eq:CDR-BRST},
it leaves tangential positive-weight fields.

For $r$ affine copies, put
\[
 P_r=I_r-\frac1r\mathbf1\mathbf1^{\mathsf t},
 \qquad
 \delta_r=\left(:\sum_{i,j}(P_r)_{ij}\gamma_i b_j:\right)_{(0)}.
\]
The tangential fields are
$\bar\gamma=r^{-1}\sum\gamma_i$,
$\bar c=r^{-1}\sum c_i$,
$\bar\beta=\sum\beta_i$ and $\bar b=\sum b_i$;
these have unit pairings. Their complement has
$r-1$ normal quartets.
For a partition into blocks let $Q$ be the
orthogonal projector onto block-constant vectors,
with entries $1/|B|$ within a block $B$. Then
\begin{equation}\label{cyprim:eq:diagonal-regrouping}
 P_r=(I_r-Q)+
       \left(Q-\frac1r\mathbf1\mathbf1^{\mathsf t}\right).
\end{equation}
The summands are orthogonal projectors onto
within-block differences and differences of
block means. Their differentials add to
$\delta_r$. Further refinement gives the same
full projector. Changing normal basis sends
$v,c$ by its matrix and $\beta,b$ by its inverse
transpose, preserving OPEs and $\sum:vb:$.
Matrix composition proves compatibility of
every affine regrouping.
These are target diagonals in $Y^r$. Insertion
diagonals for supported chiral collisions have
different ambient spaces and coefficient functors;
a comparison must relate their operations.

\subsection{The global absolute normal bundle}

Let $E\to Y$ be a holomorphic vector bundle and consider the
chiral de Rham sheaf on $\operatorname{Tot}(E)$,
with fibrewise polynomial states or the same
fixed-weight formal completion. In a local frame write
$u^i,v^a$ for base and fibre coordinates and $c_u^i,c_v^a$
for their odd partners. The fibre Euler
vector field $\mathsf E=v^a\partial_{v^a}$ is
intrinsic. Its chiral contraction is intrinsic
under~\eqref{cyprim:eq:CDR-coordinate}, so
$\delta=(\iota_{\mathsf E})_{(0)}$ is global.

The global grading now uses $-J$ for every
fermion, including the base ones. Indeed a frame
change $v'=g(u)v$ gives
$c_v'=g(u)c_v+(dg)_u(v)c_u$; a normal-only grading
would not respect this transition. Intrinsic
fermion parity is absorbed into this cochain
grading. Locally the differential is the normal
quartet differential and kills the tangential
fields.

For the vertical frame vector field
$W=A^b_a(u)v^a\partial_{v^b}$, the current
\[
 H_A=:A^b_ac_v^a\beta_b:
       -:\partial_iA^b_a\,c_u^ic_v^ab_b:
\]
has $[\delta,(H_A)_{(0)}]=\mathcal L_W$, where
$\mathcal L_W$ is the infinitesimal chiral de Rham action of $W$.
Apply the odd derivation to the two terms: the
result is
$:A^b_av^a\beta_b:+:A^b_ac_v^ab_b:
+:\partial_iA^b_a\,v^ac_u^ib_b:$,
the full chiral Lie derivative. The sign of the
last term uses $\delta(c_u^ic_v^ab_b)=-c_u^iv^ab_b$.

\begin{theorem}\label{cyprim:thm:absolute-normal}
The cohomology vertex sheaf of the absolute
Euler complex is canonically
$\Omega_Y^{\mathrm{ch}}$, with zero differential
and cochain grading $-J$. The assertion holds
on both the polynomial and fixed-weight formal
carriers.
\end{theorem}

\begin{proof}
The local normal contraction leaves precisely
the tangential states. Locally any varying
frame matrix is a constant matrix times
$\exp A(u)$. The constant factor fixes the
normal vacuum and tangential states. For the
varying factor, let $F_s$ be the frame action
of $\exp(sA)$. It commutes with $\delta$, and
the current identity gives
\[
 F_1-1=\delta K+K\delta,
 \qquad K=\int_0^1(H_A)_{(0)}F_s\,ds.
\]
Hence vertical frame changes induce the identity
on local cohomology. Base coordinate changes
induce exactly the base chiral de Rham maps.
These identifications compose on triple overlaps
and preserve the vertex products.
\end{proof}

No global strict inclusion of tangential vertex
fields is asserted. For the full diagonal in
$Y^r$, its linear normal bundle is
$T_Y\otimes(\mathbb C^r/\mathbb C_{\mathrm{diag}})$;
the theorem applies to this bundle and its
linear regroupings.

\subsection{The actual formal diagonal}

The normal bundle and the formal neighbourhood
are different geometric objects. In coordinates
$(u,u+v)$, a change $u'=f(u)$ gives
\[
 v'=f(u+v)-f(u)
    =df_u(v)+\tfrac12d^2f_u(v,v)+O(v^3).
\]
A holomorphic formal linearization over the
first projection, inducing the identity on
the normal bundle, would have quadratic term
$\operatorname{Exp}_u(v)=u+v-\frac12\Gamma_u(v,v)+O(v^3)$.
Here $\Gamma_u$ is symmetric in its two inputs.
Comparing coordinates gives
\[
 \Gamma_{f(u)}(df_uv,df_uv)
     =df_u\Gamma_u(v,v)-d^2f_u(v,v).
\]
This is the transformation law of a torsion-free
holomorphic connection on $T_Y$. On a complex
surface its curvature would have type $(2,0)$,
so the degree-four Chern--Weil form for $c_2$
would have type $(4,0)$ and vanish. This contradicts
$\int_Yc_2=12$ or $24$ for Enriques or K3.
It rules out this holomorphic linearization.
The perfect diagonal kernel $\mathcal O_\Delta$
still exists globally, with local Koszul
resolutions by the coordinate differences.

For K\"ahler $Y$, a smooth formal exponential
instead identifies the Dolbeault algebra of
the formal diagonal with
\begin{equation}\label{cyprim:eq:formal-base}
 \left(\mathcal A_Y^{0,*}
       (\widehat{\operatorname{Sym}}T_Y^*),
       D_B=\bar\partial+R\right),
 \quad R=\sum_{n\ge2}R_n,\quad
 \bar\partial R+\tfrac12[R,R]=0.
\end{equation}
Here $R_n$ is a $(0,1)$-form valued in homogeneous
vertical vector fields of degree $n$. We use
the actual Taylor coefficients, including their
factorials; $R_2$ represents the Atiyah class
\cite[Theorem~3.2]{rewriteYu2015}.

Let $\mathcal V_{\mathrm{rel}}$ be the relative formal
$bc$--$\beta\gamma$ system of $T_Y$. Its coefficients
$\mathcal A_Y^{0,*}$ are central vertex elements
with translation zero; there are no base
momentum fields. Use total cochain degree
Dolbeault degree minus $J$. A vertical vector
field $V=V^i(v)\partial_{v^i}$ has the lift
\[
 \mathcal L_V=
  \left(:V^i(v)\beta_i:
      +:\partial_jV^i(v)c^jb_i:\right)_{(0)}.
\]
The bosonic and fermionic anomalies cancel,
giving $[\mathcal L_V,\mathcal L_W]
=\mathcal L_{[V,W]}$.
Consequently
\begin{equation}\label{cyprim:eq:formal-field-differential}
 D_{\mathrm{form}}=\bar\partial+\mathcal L_R
\end{equation}
is a global square-zero vertex derivation,
preserving weight and charge and restricting
to $D_B$ on functions. Indeed its square is
$\mathcal L_{\bar\partial R+[R,R]/2}$.
The full chiral coordinate transformation
intertwines it under a change of formal
exponential: the base-dependent derivative
of that transformation gives the same
conjugation term as for $\bar\partial+R$.
The coordinate composition law proves the
exact cocycle identity. Fixed weight bounds
the number of contractions that lower the
degree of a coefficient, so these operations
are continuous in the specified completion.

The Euler contraction does not commute with
this differential. Write $\iota_V=:V^ib_i:$ for chiral
contraction, extending to form-valued vector fields with
the tensor signs. The Cartan identity and
$[V_n,\mathsf E]=-(n-1)V_n$ give
\begin{equation}\label{cyprim:eq:curvature-error}
 [D_{\mathrm{form}},\delta]
       =-\sum_{n\ge2}(n-1)(\iota_{R_n})_{(0)}.
\end{equation}
Thus $D_{\mathrm{form}}+\delta$ need not square
to zero. The curved formal neighbourhood
requires more than adding the linear contraction.

\subsection{A conjugated resolution of the central base}

A different construction resolves the central
base while retaining explicit curvature terms.
For a homogeneous degree-$n$ vector field set
\[
 H_{V_n}=\frac1n:\partial_jV_n^i c^j\beta_i:,
 \qquad
 \delta H_{V_n}=K_{V_n}
 =:V_n^i\beta_i:
       +\frac1n:\partial_jV_n^ic^jb_i:.
\]
Euler's identity gives the first term; the
odd derivation sign gives the second. Its
coefficient $1/n$ distinguishes this lift
from $\mathcal L_{V_n}$.
Let $\mathscr S=\sum_nH_{R_n}$ and
$S_0=\mathscr S_{(0)}$. Form degree one and
ghost degree $-1$ make $S_0$ even of total
degree zero. Put
\begin{equation}\label{cyprim:eq:conjugated-Koszul}
 d_0=\bar\partial+\delta,\qquad
 D_K=e^{S_0}d_0e^{-S_0}.
\end{equation}
On a complex $d$-fold the exponentials are
finite since $S_0^{d+1}=0$ by form degree.

\begin{proposition}\label{cyprim:prop:curved-Koszul}
The operator $D_K$ is a global square-zero
vertex derivation. The central inclusion
$(\mathcal A_Y^{0,*},\bar\partial)
\to(\mathcal V_{\mathrm{rel}},D_K)$ is a
quasi-isomorphism. For two choices $S_0,S'_0$,
the intertwiner $e^{S'_0}e^{-S_0}$ composes
exactly with the corresponding intertwiners.
\end{proposition}

\begin{proof}
The currents $H_{R_n}$ transform tensorially
under relative linear frames with central
coefficients. Their zero modes are global
derivations, and the finite exponentials
are vertex automorphisms fixing the base.
Conjugation proves square-zero.
The normal-number homotopy is equivariant
under these frames. With its tensor sign
on forms, it anticommutes with $\bar\partial$,
so $d_0h+hd_0=1-i\pi$ globally.
Transport it to
\[
 h_K=e^{S_0}he^{-S_0},\qquad
 \pi_K=\pi e^{-S_0},\qquad
 D_Kh_K+h_KD_K=1-i\pi_K.
\]
This proves the cohomology assertion. The
intertwiner identities follow by cancelling
adjacent exponentials. Neither projection
nor homotopy is claimed to preserve vertex
products.
\end{proof}

Conjugation gives a square-zero differential for any
such homogeneous tensors. The Maurer--Cartan equation
is needed when comparing with the actual formal diagonal.
In particular $D_K$ is not identified with
$D_{\mathrm{form}}$. It resolves only the
central base Dolbeault algebra, whereas the
absolute Euler complex retains the base
chiral de Rham fields.

On a surface all corrections are visible.
For $K_1=\sum_nK_{R_n}$, the total sign gives
$d_0\mathscr S=\bar\partial\mathscr S-K_1$, hence
\begin{equation}\label{cyprim:eq:surface-curvature}
 D_K=d_0+(K_1)_{(0)}
      -(\bar\partial\mathscr S)_{(0)}
      +\tfrac12[S_0,(K_1)_{(0)}].
\end{equation}
Terms of form degree at least three vanish.
The first-form action on functions is $R$;
the second-form term can contain negative
ghost fields, so functions need not be a
subcomplex for $D_K$.

For an explicit calculation, take fibre
coordinates $v_1,v_2$ and exterior form
variables $\eta_1,\eta_2$, ordered before
the odd fibre variables. Let
$V=v_1^2\partial_{v_2}$,
$W=v_2^2\partial_{v_1}$ and
$S_0=\eta_1(:v_1c_1\beta_2:)_{(0)}
+\eta_2(:v_2c_2\beta_1:)_{(0)}$.
Suppress $\bar\partial$, write
$A_1=[S_0,\delta]$, $A_2=[S_0,A_1]/2$,
and compute by single contractions:
\[
\begin{split}
 A_1v_1&=\eta_2v_2^2,\qquad A_1v_2=\eta_1v_1^2,\\
 A_2v_1&=-\eta_1\eta_2
       \left(\tfrac32v_1v_2c_1+\tfrac12v_1^2c_2\right),\\
 A_2v_2&=\eta_1\eta_2
       \left(\tfrac12v_2^2c_1+\tfrac32v_1v_2c_2\right).
\end{split}
\]
For example $S_0v_1=\eta_2v_2c_2$,
$S_0A_1v_1=-2\eta_1\eta_2v_1v_2c_1$, and
$A_1S_0v_1=\eta_1\eta_2(v_1v_2c_1+v_1^2c_2)$.
Their half-difference gives the displayed
second correction. These are polynomial
tensors, not the curvature of a specified
K3 metric.

\subsection{Charge restoration retains torsion}

To offset the charge decrease of $\delta$,
introduce an even parameter $u$ of charge
$+1$ and cochain degree zero. Then $u\delta$
preserves combined charge. Let $K$ be a
finite normal PBW subcomplex stable under
$\delta,h$, with only its scalar vacuum
in cohomology. Write $B^k=\operatorname{im}
(\delta:K^{k-1}\to K^k)$. Since $u$ acts
injectively on $K[u]$, its cycles are
$(\ker\delta)[u]$ and its boundaries are
$uB[u]$. Splitting the vacuum from the
boundaries gives
\begin{equation}\label{cyprim:eq:charge-torsion}
 H^k(K[u],u\delta)\cong
 B^k\otimes\mathbb C[u]/(u)
 \oplus\begin{cases}\mathbb C[u]|0\rangle,&k=0,\\0,&k\ne0.
 \end{cases}
\end{equation}
The complex is termwise free, so its derived
fibre at $u=0$ is $(K,0)$. The assertion
extends coefficientwise to the fixed-weight
formal normal system. For two directions,
no derivative jets and total normal number
at most three, its dimension is
$1+8+32+88=129$, from $(1+t)^4/(1-t)^4$.
There are $64$ acyclic pairs and one vacuum,
hence $64$ torsion generators after rescaling.

For the curved construction, charge restoration
starts over $\mathbb C[u,u^{-1}]$:
\[
 e^{S_0/u}(\bar\partial+u\delta)e^{-S_0/u}
 =\bar\partial+u\delta+(K_1)_{(0)}+u^{-1}K_2+\cdots.
\]
On a surface the last term is
$K_2=-(\bar\partial\mathscr S)_{(0)}
+[S_0,(K_1)_{(0)}]/2$.
A form-degree-$r$ correction has charge
$r-1$ and coefficient $u^{1-r}$.
Regular extension on the unchanged carrier
requires the negative-power coefficients
to vanish, or to gain the corresponding
divisibility in a parameter-dependent family.
The polynomial example shows they can be
nonzero. Their value for a specified K3
curvature requires a separate calculation.

\section{The compact comparison}
\label{cyprim:sec:comparison}

The preceding constructions provide concrete
state objects and operations: incidence
creators, Frobenius joining, elliptic image
lattices, parameter actions, twisted fields
and normal resolutions. The compact critical
Hall theory of Chapter~\ref{cycompact:chapter}
has a different coefficient system. Its
common elliptic torsor, relative orbit pair,
zero object over $BE$, orientation and
square-root flag data remain part of that
system when a state is extracted.

One useful recipient for its reduced state
is an extension module. Let $H,R$ be unital associative
$\mathbb C$-algebras, let $I_0:H\to R$ be a unital
algebra map, and let $D:H\to R$ be a linear map. If
$D(ab)=I_0(a)D(b)+D(a)I_0(b)$, then
\begin{equation}\label{cyprim:eq:reduced-module}
 \rho(a)=\begin{pmatrix}I_0(a)&D(a)\\0&I_0(a)\end{pmatrix}
\end{equation}
acts on the right $R$-module $R\oplus R$.
Matrix multiplication is the derivation
identity, and $D(1)=0$ gives the unit.
The extension splits as an $(H,R)$-bimodule
exactly when $D(a)=bI_0(a)-I_0(a)b$ for
some $b\in R$: a right-linear section is
$q\mapsto(bq,q)$. For a finite product,
\[
 D(a_1\cdots a_n)=\sum_j
 I_0(a_1)\cdots I_0(a_{j-1})D(a_j)
 I_0(a_{j+1})\cdots I_0(a_n).
\]
The order is retained even when $R$ is
noncommutative. Degree-zero chain maps
give the dg version; a strict splitting
also requires $db=0$.
An intertwiner with this extension and its
off-diagonal coefficient is additional to
the elliptic vacuum coefficient whose
constant term is one.

\begin{problem}\label{cyprim:prob:compact-action}
Construct a charge dictionary and an action
of the reduced compact critical Hall
coefficients on independently defined
geometric state complexes. Specify the
pullbacks, oriented pushforwards and
base-change maps on extension flags, over
the common elliptic torsor. Retain the
relative unit, marked vacuum and the
extension-module coefficient
\eqref{cyprim:eq:reduced-module}.
For a chiral action, construct the supported
collision maps, including failed cuts and
their higher compatibilities. For a target
diagonal transfer, retain its curved normal
fields and derived specialization. For an
age-shifted chiral McKay comparison, use
the explicit degrees in~\eqref{cyprim:eq:age-shift}.
\end{problem}

Fix supported bar objects with coderivations
$b_A,b_B$ and comparison components $F_{<n}$
satisfying all identities below arity $n$.
The obstruction in arity $n$ is the closed
residual
\[
 R_n=(b_BF_{<n}-F_{<n}b_A)_n,\qquad
 \partial_{\mathrm{Hom}}F_n=-R_n.
\]
A contraction onto cohomology, with projection $P$,
$\partial h+h\partial=1-P$ gives the
particular solution $F_n=-hR_n$ precisely
when $PR_n=0$. The normal-coordinate functor
solves this equation for its own Hom
complexes. The normal quartet supplies
another explicit contraction. Neither
computes the compact supported residuals.
The proposed rank-$28$ differential on the unchanged
root presentation is excluded by
Theorem~\ref{cyprim:thm:Serre-obstruction};
an enlarged derived presentation remains
a distinct possibility. Likewise the
normal-coordinate determinant does not
identify a global geometric line with an
automorphic line. Such a comparison needs
its base map, character and boundary
normalization, to which we turn next.
